\documentclass[11pt]{article}
% =======================================================================
%   Shared style for CS 300 notes.
%   Usage: place `% =======================================================================
%   Shared style for CS 300 notes.
%   Usage: place `% =======================================================================
%   Shared style for CS 300 notes.
%   Usage: place `\input{../../notes-style.tex}` right after \documentclass
%   in each chapter's lectures.tex and notes.tex.
% =======================================================================

% ---- Layout -----------------------------------------------------------
\usepackage[margin=1in]{geometry}
\usepackage{parskip}        % space between paragraphs, no indent
\usepackage{microtype}      % subtle kerning/spacing improvements

% ---- Colors (muted palette, easy on light-sensitive eyes) -------------
\usepackage{xcolor}
\definecolor{pagebg}       {HTML}{FAF6F0}  % warm cream background
\definecolor{bodyfg}       {HTML}{3D3A36}  % warm dark gray body text
\definecolor{heading}      {HTML}{3D6A6B}  % muted teal
\definecolor{subheading}   {HTML}{5B7A9C}  % dusty blue
\definecolor{accent}       {HTML}{8A6B4A}  % warm tan
\definecolor{linkcol}      {HTML}{5B7A9C}  % same as subheading
\definecolor{mutedtext}    {HTML}{6B6560}  % secondary / caption text
\definecolor{rulecol}      {HTML}{C8BFB0}  % separator / rule

% Callout box palette
\definecolor{defnFrame}    {HTML}{9DB89B}  \definecolor{defnBack}    {HTML}{EEF2E8}
\definecolor{ideaFrame}    {HTML}{9FB4CD}  \definecolor{ideaBack}    {HTML}{E8EEF5}
\definecolor{warnFrame}    {HTML}{CFAE87}  \definecolor{warnBack}    {HTML}{F5EDDF}
\definecolor{noteFrame}    {HTML}{BFA6AE}  \definecolor{noteBack}    {HTML}{F2E8EE}
\definecolor{exFrame}      {HTML}{9FBDBD}  \definecolor{exBack}      {HTML}{E8F0F0}
\definecolor{codeBack}     {HTML}{F3EDE2}
\definecolor{codeBorder}   {HTML}{D4CBBA}
\definecolor{codeKeyword}  {HTML}{6B4F7A}
\definecolor{codeString}   {HTML}{8A6B4A}
\definecolor{codeComment}  {HTML}{8A857F}

% ---- Page background + base text color --------------------------------
\usepackage{pagecolor}
\pagecolor{pagebg}
\color{bodyfg}

% ---- Fonts ------------------------------------------------------------
\usepackage[T1]{fontenc}
\IfFileExists{libertine.sty}{%
    \usepackage{libertine}      % warm, readable serif
    \usepackage[scaled=0.95]{inconsolata}  % softer monospace
}{}

% ---- Hyperref ---------------------------------------------------------
\usepackage{hyperref}
\hypersetup{
    colorlinks=true,
    linkcolor=linkcol,
    urlcolor=linkcol,
    citecolor=linkcol,
    pdfborder={0 0 0},
}

% ---- Section styling --------------------------------------------------
\usepackage[explicit]{titlesec}
\titleformat{\section}
    {\color{heading}\Large\bfseries}
    {\color{heading}\thesection\hspace{0.6em}}{0pt}{#1}
\titleformat{\subsection}
    {\color{subheading}\large\bfseries}
    {\color{subheading}\thesubsection\hspace{0.5em}}{0pt}{#1}
\titleformat{\subsubsection}
    {\color{accent}\normalsize\bfseries}
    {\color{accent}\thesubsubsection\hspace{0.5em}}{0pt}{#1}
\titleformat{\paragraph}[runin]
    {\color{accent}\bfseries}{}{0pt}{#1.}

\titlespacing*{\section}{0pt}{1.6em}{0.6em}
\titlespacing*{\subsection}{0pt}{1.2em}{0.4em}

% ---- Enumerate / itemize spacing --------------------------------------
\usepackage{enumitem}
\setlist[itemize]{leftmargin=1.4em, itemsep=2pt, topsep=4pt}
\setlist[enumerate]{leftmargin=1.6em, itemsep=2pt, topsep=4pt}

% ---- Code listings ----------------------------------------------------
\usepackage{listings}
\lstdefinestyle{notescode}{
    backgroundcolor=\color{codeBack},
    basicstyle=\ttfamily\small\color{bodyfg},
    keywordstyle=\color{codeKeyword}\bfseries,
    stringstyle=\color{codeString},
    commentstyle=\color{codeComment}\itshape,
    frame=single,
    framesep=6pt,
    rulecolor=\color{codeBorder},
    showstringspaces=false,
    breaklines=true,
    xleftmargin=10pt,
    xrightmargin=10pt,
    aboveskip=8pt,
    belowskip=8pt,
}
\lstset{style=notescode}

% ---- Callout boxes (tcolorbox) ----------------------------------------
\usepackage[most]{tcolorbox}

\newtcolorbox{defnbox}[1][]{
    enhanced, breakable,
    colback=defnBack, colframe=defnFrame, coltitle=bodyfg,
    title={\bfseries Definition \textemdash\ #1},
    fonttitle=\bfseries,
    boxrule=0.5pt, arc=2pt, left=8pt, right=8pt, top=6pt, bottom=6pt,
    attach boxed title to top left={xshift=10pt, yshift=-6pt},
    boxed title style={colback=defnFrame, colframe=defnFrame, arc=2pt},
    before skip=8pt, after skip=8pt,
}
\newtcolorbox{ideabox}[1][]{
    enhanced, breakable,
    colback=ideaBack, colframe=ideaFrame, coltitle=bodyfg,
    title={\bfseries Key idea #1},
    boxrule=0.5pt, arc=2pt, left=8pt, right=8pt, top=6pt, bottom=6pt,
    attach boxed title to top left={xshift=10pt, yshift=-6pt},
    boxed title style={colback=ideaFrame, colframe=ideaFrame, arc=2pt},
    before skip=8pt, after skip=8pt,
}
\newtcolorbox{warnbox}[1][]{
    enhanced, breakable,
    colback=warnBack, colframe=warnFrame, coltitle=bodyfg,
    title={\bfseries Gotcha #1},
    boxrule=0.5pt, arc=2pt, left=8pt, right=8pt, top=6pt, bottom=6pt,
    attach boxed title to top left={xshift=10pt, yshift=-6pt},
    boxed title style={colback=warnFrame, colframe=warnFrame, arc=2pt},
    before skip=8pt, after skip=8pt,
}
\newtcolorbox{notebox}[1][]{
    enhanced, breakable,
    colback=noteBack, colframe=noteFrame, coltitle=bodyfg,
    title={\bfseries Aside #1},
    boxrule=0.5pt, arc=2pt, left=8pt, right=8pt, top=6pt, bottom=6pt,
    attach boxed title to top left={xshift=10pt, yshift=-6pt},
    boxed title style={colback=noteFrame, colframe=noteFrame, arc=2pt},
    before skip=8pt, after skip=8pt,
}
\newtcolorbox{examplebox}[1][]{
    enhanced, breakable,
    colback=exBack, colframe=exFrame, coltitle=bodyfg,
    title={\bfseries Example #1},
    boxrule=0.5pt, arc=2pt, left=8pt, right=8pt, top=6pt, bottom=6pt,
    attach boxed title to top left={xshift=10pt, yshift=-6pt},
    boxed title style={colback=exFrame, colframe=exFrame, arc=2pt},
    before skip=8pt, after skip=8pt,
}

% ---- TikZ graph styles (added 2026-04-25 for ch_10 augmentation) ------
% tcolorbox[most] already loads tikz; this block declares the libraries
% + shared `vertex` / `edge` styles used by ch_10's general directed-graph
% diagrams. ch_9's tree-style TikZ uses inline overrides and is
% unaffected. Simple circle nodes + arrow edges only -- no production
% styling.
\usetikzlibrary{arrows.meta, positioning, calc}
\tikzset{
    vertex/.style={circle, draw, minimum size=7mm, inner sep=0pt,
                   font=\small, thick},
    vsource/.style={vertex, fill=ideaBack},
    vdone/.style={vertex, fill=defnBack},
    edge/.style={->, >=Stealth, thick},
    uedge/.style={-, thick},
    treeedge/.style={->, >=Stealth, thick},
    backedge/.style={->, >=Stealth, thick, dashed, red!70!black},
    fwdedge/.style={->, >=Stealth, thick, blue!60!black},
    crossedge/.style={->, >=Stealth, thick, green!50!black, dotted},
    mstedge/.style={-, very thick, accent},
}

% ---- Title block ------------------------------------------------------
\usepackage{titling}
\pretitle{\begin{center}\color{heading}\LARGE\bfseries}
\posttitle{\end{center}\vskip -0.4em}
\preauthor{\begin{center}\color{mutedtext}\normalsize}
\postauthor{\end{center}}
\predate{\begin{center}\color{mutedtext}\small}
\postdate{\end{center}\vspace{0.4em}{\color{rulecol}\hrule height 0.4pt}\vspace{1.2em}}
` right after \documentclass
%   in each chapter's lectures.tex and notes.tex.
% =======================================================================

% ---- Layout -----------------------------------------------------------
\usepackage[margin=1in]{geometry}
\usepackage{parskip}        % space between paragraphs, no indent
\usepackage{microtype}      % subtle kerning/spacing improvements

% ---- Colors (muted palette, easy on light-sensitive eyes) -------------
\usepackage{xcolor}
\definecolor{pagebg}       {HTML}{FAF6F0}  % warm cream background
\definecolor{bodyfg}       {HTML}{3D3A36}  % warm dark gray body text
\definecolor{heading}      {HTML}{3D6A6B}  % muted teal
\definecolor{subheading}   {HTML}{5B7A9C}  % dusty blue
\definecolor{accent}       {HTML}{8A6B4A}  % warm tan
\definecolor{linkcol}      {HTML}{5B7A9C}  % same as subheading
\definecolor{mutedtext}    {HTML}{6B6560}  % secondary / caption text
\definecolor{rulecol}      {HTML}{C8BFB0}  % separator / rule

% Callout box palette
\definecolor{defnFrame}    {HTML}{9DB89B}  \definecolor{defnBack}    {HTML}{EEF2E8}
\definecolor{ideaFrame}    {HTML}{9FB4CD}  \definecolor{ideaBack}    {HTML}{E8EEF5}
\definecolor{warnFrame}    {HTML}{CFAE87}  \definecolor{warnBack}    {HTML}{F5EDDF}
\definecolor{noteFrame}    {HTML}{BFA6AE}  \definecolor{noteBack}    {HTML}{F2E8EE}
\definecolor{exFrame}      {HTML}{9FBDBD}  \definecolor{exBack}      {HTML}{E8F0F0}
\definecolor{codeBack}     {HTML}{F3EDE2}
\definecolor{codeBorder}   {HTML}{D4CBBA}
\definecolor{codeKeyword}  {HTML}{6B4F7A}
\definecolor{codeString}   {HTML}{8A6B4A}
\definecolor{codeComment}  {HTML}{8A857F}

% ---- Page background + base text color --------------------------------
\usepackage{pagecolor}
\pagecolor{pagebg}
\color{bodyfg}

% ---- Fonts ------------------------------------------------------------
\usepackage[T1]{fontenc}
\IfFileExists{libertine.sty}{%
    \usepackage{libertine}      % warm, readable serif
    \usepackage[scaled=0.95]{inconsolata}  % softer monospace
}{}

% ---- Hyperref ---------------------------------------------------------
\usepackage{hyperref}
\hypersetup{
    colorlinks=true,
    linkcolor=linkcol,
    urlcolor=linkcol,
    citecolor=linkcol,
    pdfborder={0 0 0},
}

% ---- Section styling --------------------------------------------------
\usepackage[explicit]{titlesec}
\titleformat{\section}
    {\color{heading}\Large\bfseries}
    {\color{heading}\thesection\hspace{0.6em}}{0pt}{#1}
\titleformat{\subsection}
    {\color{subheading}\large\bfseries}
    {\color{subheading}\thesubsection\hspace{0.5em}}{0pt}{#1}
\titleformat{\subsubsection}
    {\color{accent}\normalsize\bfseries}
    {\color{accent}\thesubsubsection\hspace{0.5em}}{0pt}{#1}
\titleformat{\paragraph}[runin]
    {\color{accent}\bfseries}{}{0pt}{#1.}

\titlespacing*{\section}{0pt}{1.6em}{0.6em}
\titlespacing*{\subsection}{0pt}{1.2em}{0.4em}

% ---- Enumerate / itemize spacing --------------------------------------
\usepackage{enumitem}
\setlist[itemize]{leftmargin=1.4em, itemsep=2pt, topsep=4pt}
\setlist[enumerate]{leftmargin=1.6em, itemsep=2pt, topsep=4pt}

% ---- Code listings ----------------------------------------------------
\usepackage{listings}
\lstdefinestyle{notescode}{
    backgroundcolor=\color{codeBack},
    basicstyle=\ttfamily\small\color{bodyfg},
    keywordstyle=\color{codeKeyword}\bfseries,
    stringstyle=\color{codeString},
    commentstyle=\color{codeComment}\itshape,
    frame=single,
    framesep=6pt,
    rulecolor=\color{codeBorder},
    showstringspaces=false,
    breaklines=true,
    xleftmargin=10pt,
    xrightmargin=10pt,
    aboveskip=8pt,
    belowskip=8pt,
}
\lstset{style=notescode}

% ---- Callout boxes (tcolorbox) ----------------------------------------
\usepackage[most]{tcolorbox}

\newtcolorbox{defnbox}[1][]{
    enhanced, breakable,
    colback=defnBack, colframe=defnFrame, coltitle=bodyfg,
    title={\bfseries Definition \textemdash\ #1},
    fonttitle=\bfseries,
    boxrule=0.5pt, arc=2pt, left=8pt, right=8pt, top=6pt, bottom=6pt,
    attach boxed title to top left={xshift=10pt, yshift=-6pt},
    boxed title style={colback=defnFrame, colframe=defnFrame, arc=2pt},
    before skip=8pt, after skip=8pt,
}
\newtcolorbox{ideabox}[1][]{
    enhanced, breakable,
    colback=ideaBack, colframe=ideaFrame, coltitle=bodyfg,
    title={\bfseries Key idea #1},
    boxrule=0.5pt, arc=2pt, left=8pt, right=8pt, top=6pt, bottom=6pt,
    attach boxed title to top left={xshift=10pt, yshift=-6pt},
    boxed title style={colback=ideaFrame, colframe=ideaFrame, arc=2pt},
    before skip=8pt, after skip=8pt,
}
\newtcolorbox{warnbox}[1][]{
    enhanced, breakable,
    colback=warnBack, colframe=warnFrame, coltitle=bodyfg,
    title={\bfseries Gotcha #1},
    boxrule=0.5pt, arc=2pt, left=8pt, right=8pt, top=6pt, bottom=6pt,
    attach boxed title to top left={xshift=10pt, yshift=-6pt},
    boxed title style={colback=warnFrame, colframe=warnFrame, arc=2pt},
    before skip=8pt, after skip=8pt,
}
\newtcolorbox{notebox}[1][]{
    enhanced, breakable,
    colback=noteBack, colframe=noteFrame, coltitle=bodyfg,
    title={\bfseries Aside #1},
    boxrule=0.5pt, arc=2pt, left=8pt, right=8pt, top=6pt, bottom=6pt,
    attach boxed title to top left={xshift=10pt, yshift=-6pt},
    boxed title style={colback=noteFrame, colframe=noteFrame, arc=2pt},
    before skip=8pt, after skip=8pt,
}
\newtcolorbox{examplebox}[1][]{
    enhanced, breakable,
    colback=exBack, colframe=exFrame, coltitle=bodyfg,
    title={\bfseries Example #1},
    boxrule=0.5pt, arc=2pt, left=8pt, right=8pt, top=6pt, bottom=6pt,
    attach boxed title to top left={xshift=10pt, yshift=-6pt},
    boxed title style={colback=exFrame, colframe=exFrame, arc=2pt},
    before skip=8pt, after skip=8pt,
}

% ---- TikZ graph styles (added 2026-04-25 for ch_10 augmentation) ------
% tcolorbox[most] already loads tikz; this block declares the libraries
% + shared `vertex` / `edge` styles used by ch_10's general directed-graph
% diagrams. ch_9's tree-style TikZ uses inline overrides and is
% unaffected. Simple circle nodes + arrow edges only -- no production
% styling.
\usetikzlibrary{arrows.meta, positioning, calc}
\tikzset{
    vertex/.style={circle, draw, minimum size=7mm, inner sep=0pt,
                   font=\small, thick},
    vsource/.style={vertex, fill=ideaBack},
    vdone/.style={vertex, fill=defnBack},
    edge/.style={->, >=Stealth, thick},
    uedge/.style={-, thick},
    treeedge/.style={->, >=Stealth, thick},
    backedge/.style={->, >=Stealth, thick, dashed, red!70!black},
    fwdedge/.style={->, >=Stealth, thick, blue!60!black},
    crossedge/.style={->, >=Stealth, thick, green!50!black, dotted},
    mstedge/.style={-, very thick, accent},
}

% ---- Title block ------------------------------------------------------
\usepackage{titling}
\pretitle{\begin{center}\color{heading}\LARGE\bfseries}
\posttitle{\end{center}\vskip -0.4em}
\preauthor{\begin{center}\color{mutedtext}\normalsize}
\postauthor{\end{center}}
\predate{\begin{center}\color{mutedtext}\small}
\postdate{\end{center}\vspace{0.4em}{\color{rulecol}\hrule height 0.4pt}\vspace{1.2em}}
` right after \documentclass
%   in each chapter's lectures.tex and notes.tex.
% =======================================================================

% ---- Layout -----------------------------------------------------------
\usepackage[margin=1in]{geometry}
\usepackage{parskip}        % space between paragraphs, no indent
\usepackage{microtype}      % subtle kerning/spacing improvements

% ---- Colors (muted palette, easy on light-sensitive eyes) -------------
\usepackage{xcolor}
\definecolor{pagebg}       {HTML}{FAF6F0}  % warm cream background
\definecolor{bodyfg}       {HTML}{3D3A36}  % warm dark gray body text
\definecolor{heading}      {HTML}{3D6A6B}  % muted teal
\definecolor{subheading}   {HTML}{5B7A9C}  % dusty blue
\definecolor{accent}       {HTML}{8A6B4A}  % warm tan
\definecolor{linkcol}      {HTML}{5B7A9C}  % same as subheading
\definecolor{mutedtext}    {HTML}{6B6560}  % secondary / caption text
\definecolor{rulecol}      {HTML}{C8BFB0}  % separator / rule

% Callout box palette
\definecolor{defnFrame}    {HTML}{9DB89B}  \definecolor{defnBack}    {HTML}{EEF2E8}
\definecolor{ideaFrame}    {HTML}{9FB4CD}  \definecolor{ideaBack}    {HTML}{E8EEF5}
\definecolor{warnFrame}    {HTML}{CFAE87}  \definecolor{warnBack}    {HTML}{F5EDDF}
\definecolor{noteFrame}    {HTML}{BFA6AE}  \definecolor{noteBack}    {HTML}{F2E8EE}
\definecolor{exFrame}      {HTML}{9FBDBD}  \definecolor{exBack}      {HTML}{E8F0F0}
\definecolor{codeBack}     {HTML}{F3EDE2}
\definecolor{codeBorder}   {HTML}{D4CBBA}
\definecolor{codeKeyword}  {HTML}{6B4F7A}
\definecolor{codeString}   {HTML}{8A6B4A}
\definecolor{codeComment}  {HTML}{8A857F}

% ---- Page background + base text color --------------------------------
\usepackage{pagecolor}
\pagecolor{pagebg}
\color{bodyfg}

% ---- Fonts ------------------------------------------------------------
\usepackage[T1]{fontenc}
\IfFileExists{libertine.sty}{%
    \usepackage{libertine}      % warm, readable serif
    \usepackage[scaled=0.95]{inconsolata}  % softer monospace
}{}

% ---- Hyperref ---------------------------------------------------------
\usepackage{hyperref}
\hypersetup{
    colorlinks=true,
    linkcolor=linkcol,
    urlcolor=linkcol,
    citecolor=linkcol,
    pdfborder={0 0 0},
}

% ---- Section styling --------------------------------------------------
\usepackage[explicit]{titlesec}
\titleformat{\section}
    {\color{heading}\Large\bfseries}
    {\color{heading}\thesection\hspace{0.6em}}{0pt}{#1}
\titleformat{\subsection}
    {\color{subheading}\large\bfseries}
    {\color{subheading}\thesubsection\hspace{0.5em}}{0pt}{#1}
\titleformat{\subsubsection}
    {\color{accent}\normalsize\bfseries}
    {\color{accent}\thesubsubsection\hspace{0.5em}}{0pt}{#1}
\titleformat{\paragraph}[runin]
    {\color{accent}\bfseries}{}{0pt}{#1.}

\titlespacing*{\section}{0pt}{1.6em}{0.6em}
\titlespacing*{\subsection}{0pt}{1.2em}{0.4em}

% ---- Enumerate / itemize spacing --------------------------------------
\usepackage{enumitem}
\setlist[itemize]{leftmargin=1.4em, itemsep=2pt, topsep=4pt}
\setlist[enumerate]{leftmargin=1.6em, itemsep=2pt, topsep=4pt}

% ---- Code listings ----------------------------------------------------
\usepackage{listings}
\lstdefinestyle{notescode}{
    backgroundcolor=\color{codeBack},
    basicstyle=\ttfamily\small\color{bodyfg},
    keywordstyle=\color{codeKeyword}\bfseries,
    stringstyle=\color{codeString},
    commentstyle=\color{codeComment}\itshape,
    frame=single,
    framesep=6pt,
    rulecolor=\color{codeBorder},
    showstringspaces=false,
    breaklines=true,
    xleftmargin=10pt,
    xrightmargin=10pt,
    aboveskip=8pt,
    belowskip=8pt,
}
\lstset{style=notescode}

% ---- Callout boxes (tcolorbox) ----------------------------------------
\usepackage[most]{tcolorbox}

\newtcolorbox{defnbox}[1][]{
    enhanced, breakable,
    colback=defnBack, colframe=defnFrame, coltitle=bodyfg,
    title={\bfseries Definition \textemdash\ #1},
    fonttitle=\bfseries,
    boxrule=0.5pt, arc=2pt, left=8pt, right=8pt, top=6pt, bottom=6pt,
    attach boxed title to top left={xshift=10pt, yshift=-6pt},
    boxed title style={colback=defnFrame, colframe=defnFrame, arc=2pt},
    before skip=8pt, after skip=8pt,
}
\newtcolorbox{ideabox}[1][]{
    enhanced, breakable,
    colback=ideaBack, colframe=ideaFrame, coltitle=bodyfg,
    title={\bfseries Key idea #1},
    boxrule=0.5pt, arc=2pt, left=8pt, right=8pt, top=6pt, bottom=6pt,
    attach boxed title to top left={xshift=10pt, yshift=-6pt},
    boxed title style={colback=ideaFrame, colframe=ideaFrame, arc=2pt},
    before skip=8pt, after skip=8pt,
}
\newtcolorbox{warnbox}[1][]{
    enhanced, breakable,
    colback=warnBack, colframe=warnFrame, coltitle=bodyfg,
    title={\bfseries Gotcha #1},
    boxrule=0.5pt, arc=2pt, left=8pt, right=8pt, top=6pt, bottom=6pt,
    attach boxed title to top left={xshift=10pt, yshift=-6pt},
    boxed title style={colback=warnFrame, colframe=warnFrame, arc=2pt},
    before skip=8pt, after skip=8pt,
}
\newtcolorbox{notebox}[1][]{
    enhanced, breakable,
    colback=noteBack, colframe=noteFrame, coltitle=bodyfg,
    title={\bfseries Aside #1},
    boxrule=0.5pt, arc=2pt, left=8pt, right=8pt, top=6pt, bottom=6pt,
    attach boxed title to top left={xshift=10pt, yshift=-6pt},
    boxed title style={colback=noteFrame, colframe=noteFrame, arc=2pt},
    before skip=8pt, after skip=8pt,
}
\newtcolorbox{examplebox}[1][]{
    enhanced, breakable,
    colback=exBack, colframe=exFrame, coltitle=bodyfg,
    title={\bfseries Example #1},
    boxrule=0.5pt, arc=2pt, left=8pt, right=8pt, top=6pt, bottom=6pt,
    attach boxed title to top left={xshift=10pt, yshift=-6pt},
    boxed title style={colback=exFrame, colframe=exFrame, arc=2pt},
    before skip=8pt, after skip=8pt,
}

% ---- TikZ graph styles (added 2026-04-25 for ch_10 augmentation) ------
% tcolorbox[most] already loads tikz; this block declares the libraries
% + shared `vertex` / `edge` styles used by ch_10's general directed-graph
% diagrams. ch_9's tree-style TikZ uses inline overrides and is
% unaffected. Simple circle nodes + arrow edges only -- no production
% styling.
\usetikzlibrary{arrows.meta, positioning, calc}
\tikzset{
    vertex/.style={circle, draw, minimum size=7mm, inner sep=0pt,
                   font=\small, thick},
    vsource/.style={vertex, fill=ideaBack},
    vdone/.style={vertex, fill=defnBack},
    edge/.style={->, >=Stealth, thick},
    uedge/.style={-, thick},
    treeedge/.style={->, >=Stealth, thick},
    backedge/.style={->, >=Stealth, thick, dashed, red!70!black},
    fwdedge/.style={->, >=Stealth, thick, blue!60!black},
    crossedge/.style={->, >=Stealth, thick, green!50!black, dotted},
    mstedge/.style={-, very thick, accent},
}

% ---- Title block ------------------------------------------------------
\usepackage{titling}
\pretitle{\begin{center}\color{heading}\LARGE\bfseries}
\posttitle{\end{center}\vskip -0.4em}
\preauthor{\begin{center}\color{mutedtext}\normalsize}
\postauthor{\end{center}}
\predate{\begin{center}\color{mutedtext}\small}
\postdate{\end{center}\vspace{0.4em}{\color{rulecol}\hrule height 0.4pt}\vspace{1.2em}}


\title{CS 300 -- Chapter 1 Lectures\\\large C++ Refresher: Arrays, Vectors, and Strings}
\author{Jose de Lima}
\date{\today}

\begin{document}
\maketitle

\begin{ideabox}[Chapter map]
\textbf{Where this sits in CS 300.} Chapter 1 is the \emph{runway}. Every
later chapter (linked lists, stacks/queues, trees, hash tables, heaps, graphs)
assumes you already reach for \texttt{std::vector} without thinking, know
when indexing is O(1) and when it isn't, and can trace a loop by hand. If
any of that is shaky, the rest of the course compounds the confusion.

\textbf{What you add to your toolkit here:}
\begin{itemize}
    \item A sharp mental model of \textbf{contiguous storage} -- arrays and
          \texttt{std::vector} -- including what's fast (back-end ops, random
          access) and what's slow (insert/erase in the middle).
    \item A library of \textbf{iteration patterns}: running best, counter,
          in-place transform, paired scan -- these come back in every later
          chapter disguised as ``traverse the list/tree/graph.''
    \item \textbf{Value semantics} in C++: copying, comparing, passing by
          reference, and the cost of accidentally making a deep copy.
    \item \textbf{C-string survival} so you can read legacy code without
          panicking, plus the \texttt{std::string} / \texttt{cctype} upgrades.
\end{itemize}

\textbf{Mastery checklist} -- you should be able to, without references:
\begin{enumerate}
    \item Write a \texttt{for} loop that scans a \texttt{vector<int>} keeping a
          running max \emph{and} the index where it occurred, in one pass.
    \item State the Big-O of \texttt{push\_back}, \texttt{pop\_back},
          \texttt{insert(begin(), x)}, \texttt{erase(it)}, and \texttt{at(i)},
          and say \emph{why}.
    \item Explain why \texttt{resize(n)} and a loop of \texttt{push\_back}
          produce the same end state but differ in initial values and in
          performance.
    \item Swap two elements of a vector correctly with and without
          \texttt{std::swap}, and point out why \texttt{a = b;\ b = a;} fails.
    \item Convert a small parallel-vectors design into a
          \texttt{vector<Struct>} and explain the tradeoff.
    \item Reverse a \texttt{vector} in place correctly on the first try
          (loop bound $n/2$, swap by temp).
    \item Read a C-string character-by-character, classify with
          \texttt{cctype}, and explain why \texttt{strcmp} returns an
          \texttt{int}, not a \texttt{bool}.
\end{enumerate}

\textbf{Looking ahead.} Chapter 2 formalizes \emph{problem-solving}
strategies (recursion, greedy, DP); chapter 3 formalizes the cost analysis
(Big-O) that you're using intuitively here; chapter 4 introduces the
linked-list alternative to vectors, and every performance comparison in that
chapter is a comparison \emph{against} what you learned here.
\end{ideabox}

\section{1.1 Arrays and Vectors (general concept)}

\subsection{The motivation: one name, many values}

A regular variable holds one value. That is fine when you have a handful
of distinct things to track (a wage and a salary), but it falls apart
the moment you need a \emph{list}: the four quarter scores of a
basketball game, 365 daily temperature readings, every bid in an
auction. Nobody wants to declare \texttt{score1}, \texttt{score2},
\texttt{score3}, \texttt{score4} -- and declaring
\texttt{temperature\_001} through \texttt{temperature\_365} is absurd.

\begin{defnbox}[Array]
A variable that holds a \textbf{fixed-length, ordered sequence} of
values, all of the same type, under a single name. Each value is called
an \textbf{element}, and each element has a position in the sequence
called its \textbf{index}.
\end{defnbox}

A useful analogy: a normal variable is a truck with one cargo bay; an
array is a train whose cars are all coupled together and numbered, but
which you still call by one name.

\subsection{Indexing}

The essential feature of an array is \textbf{direct access}: given an
index, you can read or write the element at that position in constant
time, without scanning from the front. You write this with bracket
notation in most languages:
\begin{lstlisting}
myArray[2]       // element at index 2
\end{lstlisting}

\begin{ideabox}[Why direct access is $O(1)$]
Array elements are stored \emph{contiguously} in memory. Given the
array's starting address and an element's size, the address of element
$i$ is simply
\[
\text{start} + i \times \text{element\_size}.
\]
No matter how long the array is, that calculation takes the same
amount of time. This is exactly why arrays feel fast, and why random
access is the benchmark other data structures are measured against.
\end{ideabox}

\begin{notebox}[The model has a name: Word-RAM]
The picture used above -- memory is an addressable sequence of fixed-size
\emph{words}; the processor reads or writes any one word in $O(1)$ -- has
a formal name in algorithm analysis: the \textbf{Word-RAM} model.
Chapter~3 adopts this language when it makes Big-O analysis precise. For
now you only need the conclusion: indexed access into contiguous memory
is one operation, no matter how big the array.
\end{notebox}

\subsection{Zero-based indexing}

In C, C++, Java, Python, and most languages you will meet, the first
element is at index \textbf{0}, not 1. An array with 4 elements has
valid indices \texttt{0, 1, 2, 3} -- \emph{not} \texttt{1, 2, 3, 4}.

This is not a pointless quirk. The index is the \emph{offset} from the
start of the array. Element 0 is 0 slots from the start; element 3 is 3
slots from the start. The math in the box above works out cleanly
because of it.

\begin{warnbox}[Off-by-one errors]
The most common bug in any language with arrays: treating the length as
the last valid index. If \texttt{scoresList} has 10 elements, valid
indices are \texttt{0}--\texttt{9}. Reading \texttt{scoresList[10]} is
a bug.
In C++ this is \textbf{undefined behaviour} -- no exception, no error
at compile time, the program just silently reads whatever memory sits
past the array. Sometimes it crashes; sometimes it ``works'' and
corrupts a neighbouring variable. Get in the habit of writing
\texttt{i < length}, never \texttt{i <= length}.
\end{warnbox}

\subsection{Arrays vs.\ vectors}

Some languages distinguish between \textbf{arrays} (fixed size, set at
creation) and \textbf{vectors} (resizable, can grow or shrink). Coral
uses the word \emph{array} loosely; C++ gives you both:

\begin{examplebox}[C++ sequence containers]
\begin{tabular}{lll}
\textbf{Type} & \textbf{Size} & \textbf{Access} \\
\hline
C-style array \texttt{int a[10]} & fixed at compile time & \texttt{a[i]} \\
\texttt{std::array<int,10>} & fixed at compile time & \texttt{a[i]} or \texttt{a.at(i)} \\
\texttt{std::vector<int>} & grows at runtime & \texttt{v[i]} or \texttt{v.at(i)} \\
\end{tabular}\\[4pt]
\texttt{operator[]} does \emph{not} bounds-check; \texttt{.at(i)}
throws \texttt{std::out\_of\_range} if \texttt{i} is invalid. The CS
300 projects use \texttt{std::vector} almost exclusively because bid
lists grow dynamically as the CSV is read.
\end{examplebox}

\begin{notebox}[Vectors are not linked lists]
The word ``vector'' in C++ does \emph{not} mean a linked list. A
\texttt{std::vector} is an array that knows how to grow -- internally
it allocates a contiguous buffer and reallocates a larger one when it
runs out of space. Random access stays $O(1)$. We will meet linked
lists in Chapter 4 and they behave completely differently.
\end{notebox}

\begin{ideabox}[The big picture for 1.1]
Arrays give you \textbf{constant-time access to any element by
position}, at the cost of fixed (or managed) size. That single
property -- direct indexing over contiguous memory -- is the reason
arrays underpin almost every other data structure in this course.
Hash tables use arrays as their bucket storage; heaps use arrays to
lay out trees; even \texttt{std::vector} is ``just'' an array with
growth built in.
\end{ideabox}

\section{1.2 Vectors (in C++)}

This is where the course drops the Coral training wheels and starts
writing real C++. A \texttt{std::vector} is the C++ workhorse container
-- the one you reach for by default when you need an ordered list --
and essentially every CS 300 project uses it.

\subsection{Declaring a vector}

Two things have to happen before you can use a vector in a source file:
\begin{lstlisting}
#include <vector>     // bring in the vector type
std::vector<int> gameScores(4);   // declare a vector of 4 ints
\end{lstlisting}

The piece in angle brackets is the \textbf{element type} -- the kind of
value every element will hold. The number in parentheses is the
\textbf{initial size}.

\begin{defnbox}[Template parameter]
\texttt{vector} is a \emph{class template} -- a blueprint that becomes
a concrete type only when you fill in the angle brackets.
\texttt{vector<int>} is one type, \texttt{vector<std::string>} is
another, and \texttt{vector<Bid>} is another. You will see the same
syntax on \texttt{std::map}, \texttt{std::set}, and every other
standard container.
\end{defnbox}

\begin{notebox}[\texttt{using namespace std;}]
The textbook writes \texttt{using namespace std;} at the top of every
example so you can say \texttt{vector} instead of \texttt{std::vector}.
It is fine for learning, but in any non-trivial file you should write
\texttt{std::vector} explicitly. Pulling the whole \texttt{std}
namespace in globally is a well-known footgun -- it can silently
collide with identifiers in your own code or in other libraries. The
projects in this course are small enough that either convention works,
but get used to typing \texttt{std::} now.
\end{notebox}

\subsection{Accessing elements}

You have two ways to read or write an element:
\begin{lstlisting}
v.at(i)   // bounds-checked: throws std::out_of_range if i is invalid
v[i]      // not bounds-checked: undefined behaviour if i is invalid
\end{lstlisting}

\begin{warnbox}[\texttt{.at()} vs \texttt{[]}]
\texttt{.at()} is the safer form and what the textbook teaches first -- if
\texttt{i} is out of range you get an exception you can actually catch.
\texttt{[]} trades safety for speed (no check on every access), which
matters in tight inner loops but not in 99\% of coursework code. Rule
of thumb: \texttt{.at()} while you are learning and when the index
came from the user; \texttt{[]} inside loops where you already
guaranteed the bound with \texttt{i < v.size()}.
\end{warnbox}

Indices start at \textbf{0}, same as arrays. A vector with 4 elements
has valid indices \texttt{0, 1, 2, 3}. The index can be any non-negative
integer expression -- a literal, a variable, the result of arithmetic.
It cannot be a floating-point value (\texttt{v.at(1.0)} will not
compile).

\begin{examplebox}[Nth-oldest lookup]
\begin{lstlisting}
std::vector<int> oldestPeople = {122, 119, 117, 117, 116};
int n;                          // user's 1-based pick
std::cin >> n;
if (n >= 1 && n <= 5) {
    std::cout << oldestPeople.at(n - 1);  // translate 1-based -> 0-based
}
\end{lstlisting}
Note the \texttt{n - 1}: the user thinks ``the 1st oldest person,''
but internally that is index \texttt{0}. This is one of the most common
places off-by-one bugs sneak in.
\end{examplebox}

\subsection{Size, looping, and \texttt{size\_t}}

Every vector knows its current length via \texttt{v.size()}. The classic
indexed loop over every element:
\begin{lstlisting}
for (std::size_t i = 0; i < v.size(); ++i) {
    std::cout << v.at(i);
}
\end{lstlisting}

\begin{defnbox}[\texttt{std::size\_t}]
The unsigned integer type the standard library uses for sizes and
indices. \texttt{v.size()} returns \texttt{size\_t}, so your loop
counter should be \texttt{size\_t} too. Using a signed
\texttt{int} instead will draw a ``comparison between signed and
unsigned'' warning on strict compilers. The textbook uses
\texttt{unsigned int}, which works but is not the idiomatic choice.
\end{defnbox}

For read-only traversals, the modern and much cleaner form is the
\textbf{range-based for loop}:
\begin{lstlisting}
for (int value : v) {
    std::cout << value;
}
\end{lstlisting}
No index, no bounds to worry about. Use this whenever you do not
actually need \texttt{i}.

\subsection{Initialisation}

Four common patterns:
\begin{lstlisting}
std::vector<int> a(5);               // 5 elements, each 0
std::vector<int> b(5, -1);           // 5 elements, each -1
std::vector<int> c = {3, 1, 4};      // 3 elements, listed explicitly
std::vector<int> d;                  // empty; push_back later
\end{lstlisting}
Pattern \texttt{d} is by far the most common in the projects -- you
start empty and \texttt{push\_back} each bid as you read the CSV.

\begin{ideabox}[Vectors grow]
Unlike a C-style array, a vector is \emph{resizable}. You can call
\texttt{v.push\_back(x)} to append, \texttt{v.pop\_back()} to remove
the last element, or \texttt{v.resize(n)} to change its length. The
vector reallocates its internal storage as needed; you do not manage
memory by hand. This is the feature that makes \texttt{vector} the
default choice when you do not yet know how big the list will be.
\end{ideabox}

\subsection{Common errors}

\begin{itemize}
    \item \textbf{Missing \texttt{\#include <vector>}} -- the compiler
          does not know what \texttt{vector} is, and the error message
          is spectacularly unhelpful (``ISO C++ forbids declaration of
          vector with no type'').
    \item \textbf{Missing \texttt{std::} or \texttt{using namespace std;}}
          -- same error, same fix.
    \item \textbf{Out-of-range access with \texttt{[]}} -- no crash, no
          diagnostic, wrong data. If you cannot reproduce a bug,
          suspect this first.
    \item \textbf{Signed/unsigned comparison in the loop} -- warning,
          not error, but occasionally causes an infinite loop if
          \texttt{i} goes negative.
\end{itemize}

\begin{ideabox}[The big picture for 1.2]
A \texttt{std::vector} is the modern C++ answer to ``I need a list.''
It combines the $O(1)$ indexed access of an array with the flexibility
of dynamic growth, and it manages its own memory. Get fluent in
declaration, \texttt{push\_back}, \texttt{size()}, and iteration (both
indexed and range-based) -- every later data-structure chapter uses
\texttt{vector} either as a direct tool or as a reference point.
\end{ideabox}

\section{1.3 Array / Vector Iteration Drill}

The textbook presents this section as three interactive drills with no
textual content:
\begin{enumerate}
    \item find the maximum value,
    \item count the negatives,
    \item bubble the largest value to the end.
\end{enumerate}
Each drill is a specific instance of a general pattern. Writing those
patterns out once makes the drills easier -- and more importantly,
every iterative algorithm later in the course is a remix of them.

\subsection{Pattern 1 -- Linear scan with a running ``best''}

\begin{ideabox}[The pattern]
Keep a variable that holds the best value seen so far. Visit every
element; whenever the current element beats the running best, update
the running best. At the end, the variable holds the answer.
\end{ideabox}

\paragraph{Find the maximum.}
\begin{lstlisting}
int maxVal = v.at(0);                // seed with the first element
for (std::size_t i = 1; i < v.size(); ++i) {
    if (v.at(i) > maxVal) {
        maxVal = v.at(i);
    }
}
\end{lstlisting}

\begin{warnbox}[Never seed a max with 0]
A tempting but broken shortcut is \texttt{int maxVal = 0;} before the
loop. If \emph{every} value in the vector is negative, you will
incorrectly report \texttt{0} as the max. Seed with the first element
(or with the smallest possible value for the type,
\texttt{std::numeric\_limits<int>::min()}) -- not with an arbitrary
constant. The same trap applies to finding the minimum: do not seed
with \texttt{0}.
\end{warnbox}

\paragraph{Find the minimum.} Same pattern, flip the comparison:
\texttt{if (v.at(i) < minVal)}.

\paragraph{Find the \emph{index} of the maximum.} When you need the
position (e.g., for sorting), track \texttt{i} instead of the value:
\begin{lstlisting}
std::size_t maxIdx = 0;
for (std::size_t i = 1; i < v.size(); ++i) {
    if (v.at(i) > v.at(maxIdx)) {
        maxIdx = i;
    }
}
\end{lstlisting}

\subsection{Pattern 2 -- Linear scan with a running counter}

\begin{ideabox}[The pattern]
Start a counter at zero. Visit every element; whenever it matches the
condition of interest, increment the counter. At the end, the counter
holds the total.
\end{ideabox}

\paragraph{Count the negatives.}
\begin{lstlisting}
int negCount = 0;
for (std::size_t i = 0; i < v.size(); ++i) {
    if (v.at(i) < 0) {
        ++negCount;
    }
}
\end{lstlisting}

The same shape solves ``how many even numbers,'' ``how many bids above
\$100,'' ``how many strings start with \texttt{A}.'' The condition
changes; the structure does not.

\paragraph{Running sum / running product.} A tiny variation: instead
of incrementing by 1, accumulate the element itself. \texttt{sum +=
v.at(i)} gives a total; \texttt{product *= v.at(i)} gives a product
(seed with 1, not 0).

\subsection{Pattern 3 -- One selection-sort pass}

The third drill asks you to move the largest value to the last
position of the array. This is exactly one pass of \textbf{selection
sort}, and it combines Patterns 1 (find the max index) with a
\textbf{swap}:
\begin{lstlisting}
// After this loop, the largest value sits at v.back().
std::size_t maxIdx = 0;
for (std::size_t i = 1; i < v.size(); ++i) {
    if (v.at(i) > v.at(maxIdx)) {
        maxIdx = i;
    }
}
std::swap(v.at(maxIdx), v.at(v.size() - 1));
\end{lstlisting}

\begin{notebox}[Why the swap matters]
If you simply \emph{copy} \texttt{v.at(maxIdx)} into the last slot,
you overwrite whatever was there. The value that \emph{was} at the end
is gone. A \textbf{swap} preserves it by putting it into the slot the
max came from. This is the same reason selection sort is an
\emph{in-place} sort: every value stays in the array, just in a new
position.
\end{notebox}

\subsection{Range-based form}

When you do not need the index, the same three patterns read much more
cleanly:
\begin{lstlisting}
int maxVal = v.at(0);
for (int x : v) { if (x > maxVal) maxVal = x; }

int negCount = 0;
for (int x : v) { if (x < 0) ++negCount; }
\end{lstlisting}

Use the indexed form when you need \texttt{i} (e.g., for the swap in
Pattern 3) and the range-based form otherwise.

\begin{ideabox}[The big picture for 1.3]
Almost every array/vector algorithm you will write is a combination of
three building blocks: \emph{scan with a running best}, \emph{scan
with a running count or sum}, and \emph{scan-plus-swap}. Internalise
the shape of each loop -- the initialization, the condition, the
update -- and you will recognize them on sight in every later
algorithm (selection sort, linear search, minimum finding, summing a
column, etc.).
\end{ideabox}

\section{1.4 Iterating Through Vectors}

Section~1.3 covered the \emph{patterns} (running best, counter, scan-plus-swap).
This section is about the \emph{loop idiom} those patterns are built on,
and about the standard-library tools that often replace the loop entirely.

\subsection{Anatomy of the canonical loop}

The loop shape you will type thousands of times in C++:
\begin{lstlisting}
for (std::size_t i = 0; i < v.size(); ++i) {
    // work with v.at(i) or v[i]
}
\end{lstlisting}

Three pieces, each doing one job:

\begin{itemize}
    \item \textbf{Init: \texttt{std::size\_t i = 0}.}\\
          Declares the loop counter \emph{inside} the \texttt{for}
          header so it does not leak into the surrounding scope.
          \texttt{size\_t} is the type \texttt{v.size()} returns;
          matching it silences the ``signed/unsigned comparison''
          warning. Start at \texttt{0} because indices are zero-based.

    \item \textbf{Condition: \texttt{i < v.size()}.}\\
          Continue while \texttt{i} is a valid index. Valid indices
          for a vector of length $n$ are $0, 1, \ldots, n{-}1$. The
          last valid index is \texttt{v.size() - 1}, not
          \texttt{v.size()}, which is why the condition is strict
          less-than.

    \item \textbf{Update: \texttt{++i}.}\\
          Advance to the next index. Prefer \texttt{++i} to
          \texttt{i++}; see the aside below.
\end{itemize}

\begin{warnbox}[\texttt{<} vs.\ \texttt{<=}]
Writing \texttt{i <= v.size()} runs the loop one too many times and
calls \texttt{v.at(v.size())}, which throws
\texttt{std::out\_of\_range}. The error looks like:
\begin{lstlisting}
terminate called after throwing an instance of 'std::out_of_range'
  what():  vector::_M_range_check
\end{lstlisting}
Fix: use \texttt{<}, not \texttt{<=}. Every time.
\end{warnbox}

\begin{notebox}[\texttt{++i} vs.\ \texttt{i++}]
Both advance \texttt{i} by one. The difference is what the expression
\emph{evaluates to}: \texttt{++i} returns the new value (after
increment), \texttt{i++} returns a copy of the old value (before
increment). For a plain \texttt{int} the compiler optimises the
copy away and there is no cost difference -- but for richer types
(iterators, \texttt{BigInt}, etc.) \texttt{i++} may actually construct
a temporary. Building the \texttt{++i} habit means you never have to
think about it. In a \texttt{for} header, where the return value is
thrown away, \texttt{++i} is strictly the better default.
\end{notebox}

\subsection{Empty loops and infinite loops}

Two classic mistakes with the \texttt{for} header:

\begin{itemize}
    \item \textbf{Missing condition} -- \texttt{for (i = 0; ; ++i)}.
          With no condition the loop never exits on its own. It will
          eventually walk off the end of the vector and crash.
    \item \textbf{Condition that never becomes false} --
          \texttt{for (i = 0; i < 99; ++i)} over a vector of 365
          elements visits only the first 99; a condition referring to
          something that is not the vector's size is nearly always a
          bug. Use \texttt{v.size()} so the loop adapts to whatever
          the vector actually is.
\end{itemize}

\begin{warnbox}[Reverse iteration with \texttt{size\_t}]
It is tempting to iterate backwards with
\texttt{for (std::size\_t i = v.size() - 1; i >= 0; --i)}. \textbf{This
is an infinite loop.} \texttt{size\_t} is unsigned -- it cannot go
negative. After \texttt{i == 0}, \texttt{--i} wraps around to a
gigantic positive number. Safe patterns:
\begin{lstlisting}
// Loop while i > 0, then handle 0 separately:
for (std::size_t i = v.size(); i-- > 0; ) {
    // i now takes values size-1, size-2, ..., 0
}
\end{lstlisting}
or use a signed counter when the index arithmetic needs to go
negative.
\end{warnbox}

\subsection{Common errors}

\begin{itemize}
    \item \textbf{Out-of-range access} with \texttt{.at()} throws
          \texttt{std::out\_of\_range} (program aborts unless you
          catch it). With \texttt{[]} it is silent undefined
          behaviour. Bounds-check the index whenever it came from the
          user.
    \item \textbf{Seeding \texttt{maxVal} or \texttt{minVal} with 0.}
          Covered in 1.3. Seed with \texttt{v.at(0)} instead.
    \item \textbf{Mixing signed and unsigned in the comparison.}
          \texttt{int i} vs.\ \texttt{v.size()} produces warnings and
          can surprise you if \texttt{i} ever goes negative. Use
          \texttt{size\_t}.
\end{itemize}

\subsection{The standard-library alternative}

Most of the loops in 1.3 and this section have a one-liner equivalent
in \texttt{<algorithm>} and \texttt{<numeric>}. Knowing these saves a
surprising amount of noise in your code:

\begin{lstlisting}
#include <algorithm>
#include <numeric>

// Sum of all elements
int total = std::accumulate(v.begin(), v.end(), 0);

// Maximum element (returns an iterator; dereference with *)
int maxVal = *std::max_element(v.begin(), v.end());

// Count elements matching a condition
int negCount = std::count_if(v.begin(), v.end(),
                             [](int x) { return x < 0; });
\end{lstlisting}

\begin{ideabox}[Iterators, a 30-second preview]
\texttt{v.begin()} and \texttt{v.end()} are \textbf{iterators} --
lightweight objects that behave like pointers into the container.
\texttt{begin()} points at the first element; \texttt{end()} points
\emph{one past} the last (the same ``one-past-the-end'' convention as
\texttt{i < v.size()}). Almost every standard-library algorithm
takes an iterator pair. You do not need to master iterators now, but
recognize the \texttt{.begin(), .end()} pattern -- it reappears on
every STL container.
\end{ideabox}

\begin{notebox}[When to write the loop vs.\ call the algorithm]
The standard-library form is clearer \emph{when the operation has a
name} (sum, max, count, sort). The hand-written loop is clearer when
you need something more specific -- e.g., ``sum only the bids above
\$100 that belong to the closed funds.'' As a rule of thumb:
one-condition scans $\rightarrow$ algorithm; multi-condition or
index-aware work $\rightarrow$ write the loop.
\end{notebox}

\begin{ideabox}[The big picture for 1.4]
The classic indexed loop -- \texttt{for (size\_t i = 0; i < v.size();
++i)} -- is the default shape for vector iteration. Know every piece
of it cold, know the off-by-one and signed/unsigned traps, and know
that a handful of standard-library calls
(\texttt{accumulate}, \texttt{max\_element}, \texttt{count\_if}) can
replace a lot of handwritten scans when you recognize the pattern.
\end{ideabox}

\section{1.5 Multiple Vectors}

\subsection{The parallel-vectors technique}

Sometimes the thing you need to store has more than one field per
``row.'' A standard example pairs country names with the average minutes
of TV watched per day in each country:

\begin{lstlisting}
std::vector<std::string> ctryNames(5);
std::vector<int>         ctryMins(5);

ctryNames.at(0) = "China";   ctryMins.at(0) = 155;
ctryNames.at(1) = "Sweden";  ctryMins.at(1) = 154;
ctryNames.at(2) = "Russia";  ctryMins.at(2) = 246;
ctryNames.at(3) = "UK";      ctryMins.at(3) = 216;
ctryNames.at(4) = "USA";     ctryMins.at(4) = 274;
\end{lstlisting}

The convention: \textbf{same index, same row}. Index 3 means ``UK,
216 minutes,'' no matter which vector you look in. To search, you scan
one vector for the key and read the matching index in the other:

\begin{lstlisting}
for (std::size_t i = 0; i < ctryNames.size(); ++i) {
    if (ctryNames.at(i) == userCountry) {
        std::cout << ctryMins.at(i) << '\n';
        break;              // early exit once we find it
    }
}
\end{lstlisting}

\subsection{Short-circuiting the search}

The textbook uses a flag variable to stop the loop early:
\begin{lstlisting}
bool foundCountry = false;
for (i = 0; (i < ctryNames.size()) && (!foundCountry); ++i) {
    if (ctryNames.at(i) == userCountry) {
        foundCountry = true;
        /* ... */
    }
}
\end{lstlisting}

This works, but it is not the cleanest form. Two better options:

\begin{itemize}
    \item \textbf{\texttt{break}} out of the loop directly -- no flag,
          one less variable to reason about.
    \item \textbf{\texttt{std::find}} from \texttt{<algorithm>}
          expresses ``scan for a value'' in one line:
\begin{lstlisting}
auto it = std::find(ctryNames.begin(), ctryNames.end(), userCountry);
if (it != ctryNames.end()) {
    std::size_t i = it - ctryNames.begin();  // index from iterator
    std::cout << ctryMins.at(i) << '\n';
}
\end{lstlisting}
\end{itemize}

\begin{defnbox}[Linear search]
Walking a sequence one element at a time until you find the target
(or run off the end) is a \textbf{linear search}. Its worst-case cost
is $O(n)$ -- proportional to the size of the vector. For small
lists this is fine; for hundreds of thousands of bids, it is
painfully slow. A big chunk of CS 300 -- hash tables, BSTs -- exists
to beat this $O(n)$ lookup.
\end{defnbox}

\subsection{Why parallel vectors are a smell}

Parallel vectors are a reasonable starting technique, but they are
\textbf{fragile}. Two specific ways they go wrong:

\begin{itemize}
    \item \textbf{They drift out of sync.} Insert into one vector and
          forget the other; now index 3 means ``UK, 246 minutes''
          instead of ``UK, 216.'' The compiler will not warn you.
    \item \textbf{They do not scale past two or three fields.} When a
          ``row'' has ten fields you end up with ten vectors you have
          to keep in lockstep -- and every operation (sort, delete,
          filter) has to touch all of them.
\end{itemize}

\begin{warnbox}[The real answer: a vector of structs]
In practice you group related fields into a single \textbf{struct}
(or class) and keep one vector of those. The TV example rewritten the
right way:
\begin{lstlisting}
struct Country {
    std::string name;
    int         minutesTV;
};

std::vector<Country> countries = {
    {"China",  155},
    {"Sweden", 154},
    {"Russia", 246},
    {"UK",     216},
    {"USA",    274},
};

for (const Country& c : countries) {
    if (c.name == userCountry) {
        std::cout << c.minutesTV << '\n';
        break;
    }
}
\end{lstlisting}
One collection, fields that cannot drift apart, a single sort will
move names and minutes together. This is the form the CS 300 projects
use: a \texttt{Bid} struct with \texttt{bidId}, \texttt{title},
\texttt{fund}, \texttt{amount}, stored in a \texttt{std::vector<Bid>}.
Parallel vectors appear in intro material because structs have not yet
been introduced; you will almost never write them outside an intro
exercise.
\end{warnbox}

\subsection{When parallel vectors are actually okay}

Two niches:
\begin{itemize}
    \item \textbf{Temporary scratch data} -- e.g., during sorting you
          sometimes need a parallel ``order of original indices''
          vector.
    \item \textbf{Performance-critical numerical code} -- storing each
          field contiguously (``struct of arrays'' layout) can make
          SIMD and cache behaviour better. This is a microoptimisation
          you will not need in CS 300.
\end{itemize}

\begin{ideabox}[The big picture for 1.5]
The pattern to learn here is the idea of \textbf{linking data across
indices}: position $i$ in one collection corresponds to position $i$
in another. Parallel vectors are a primitive, error-prone way to do
this. The same idea reappears, better expressed, in a vector of
structs (``one row per element''), in a hash table (``key $\to$
value''), and in a pair of iterators. Recognize it early.
\end{ideabox}

\section{1.6 Vector Resize}

\subsection{Declaring without a size}

The size given in the declaration is not mandatory:
\begin{lstlisting}
std::vector<int> carSales;   // empty: size 0, no elements yet
\end{lstlisting}
At this point any access -- \texttt{carSales.at(0)},
\texttt{carSales[0]}, \texttt{carSales.front()} -- is invalid, because
there is nothing there. You have to grow the vector first.

\subsection{\texttt{resize(n)}}

\texttt{resize(n)} sets the vector's size to exactly \texttt{n}.

\begin{itemize}
    \item If \texttt{n} is \emph{larger} than the current size, new
          elements are appended at the end. They are
          \textbf{default-initialised}: \texttt{0} for numeric types,
          empty string for \texttt{std::string}, default-constructed
          for your own structs/classes.
    \item If \texttt{n} is \emph{smaller}, the extra elements at the
          end are destroyed. Anything stored past index \texttt{n-1}
          is gone -- no warning, no recycle bin.
    \item If \texttt{n} equals the current size, \texttt{resize} does
          nothing.
\end{itemize}

\begin{lstlisting}
std::vector<int> v;
v.resize(3);     // v is now {0, 0, 0}
v.at(0) = 122;
v.at(1) = 11;
v.at(2) = 7;     // v is {122, 11, 7}

v.resize(2);     // v is {122, 11}; the 7 is gone
v.at(2);         // throws std::out_of_range
\end{lstlisting}

\subsection{\texttt{resize()} vs.\ \texttt{push\_back()} -- which one?}

Two different grow-the-vector strategies, and it is worth knowing when
each fits.

\begin{examplebox}[Two ways to read $n$ user values]
\textbf{Resize first} (the textbook's preferred style):
\begin{lstlisting}
int n;
std::cin >> n;
v.resize(n);
for (std::size_t i = 0; i < v.size(); ++i) {
    std::cin >> v.at(i);
}
\end{lstlisting}

\textbf{Push-back as you go:}
\begin{lstlisting}
int n;
std::cin >> n;
for (int i = 0; i < n; ++i) {
    int x;
    std::cin >> x;
    v.push_back(x);
}
\end{lstlisting}

Both leave \texttt{v} with \texttt{n} elements. The difference is who
knows the count up front. When the count is known in advance, either
works. When the count is \emph{not} known -- ``read bids until
end-of-file'' -- \texttt{push\_back} is the only sensible choice,
because you cannot call \texttt{resize} with a number you do not
have yet.
\end{examplebox}

\begin{ideabox}[Rule of thumb]
Use \texttt{resize(n)} when you know exactly how many elements you
want and want them default-initialised in one shot. Use
\texttt{push\_back} when the count is discovered one element at a time
(the common case for reading files, parsing input, building results
on the fly). CS 300 projects use \texttt{push\_back} for the bid list.
\end{ideabox}

\subsection{Size, capacity, and \texttt{reserve()} (what intro texts omit)}

A vector tracks \emph{two} numbers, not one:

\begin{defnbox}[Size vs.\ capacity]
\begin{itemize}
    \item \textbf{size} (\texttt{v.size()}): how many elements the
          vector currently \emph{holds}. This is what your loops
          iterate over.
    \item \textbf{capacity} (\texttt{v.capacity()}): how many elements
          its internal buffer \emph{could} hold before a reallocation
          is needed. Always $\geq$ size.
\end{itemize}
\end{defnbox}

When you \texttt{push\_back} and size is about to exceed capacity, the
vector allocates a bigger internal buffer (typically $1.5\times$ or
$2\times$ the old capacity), copies everything over, and frees the
old buffer. That reallocation is $O(n)$ -- but because it happens
geometrically less often as the vector grows, the \emph{amortised}
cost of \texttt{push\_back} is $O(1)$.

\begin{defnbox}[Amortized cost]
A data-structure operation has \textbf{amortized cost} $T(n)$ if
\emph{any} sequence of $k$ operations costs at most $k \cdot T(n)$ in
total. So $\Theta(1)$ amortized doesn't mean every individual call is
cheap -- it means the expensive ones happen rarely enough that the
average over many calls is bounded. \texttt{push\_back} fits exactly:
most calls write one slot in $O(1)$; the occasional reallocation costs
$O(n)$, but a doubling policy means it can only happen once per
\emph{geometric} doubling, so the total bill stays $O(k)$ for $k$ pushes.
Chapter~3 derives this rigorously.
\end{defnbox}

If you know ahead of time roughly how big the vector will get, you
can skip the intermediate reallocations with \texttt{reserve}:
\begin{lstlisting}
std::vector<Bid> bids;
bids.reserve(10000);         // capacity -> 10000, size still 0
for (/* each CSV row */) {
    bids.push_back(parseRow(row));  // no reallocations
}
\end{lstlisting}

\begin{warnbox}[\texttt{resize} vs.\ \texttt{reserve} -- not the same]
\begin{itemize}
    \item \texttt{resize(n)} changes the \emph{size} to \texttt{n}
          (creating or destroying elements as needed). Your loops
          now iterate \texttt{n} times.
    \item \texttt{reserve(n)} only changes the \emph{capacity};
          size stays where it was. Nothing to iterate over until you
          \texttt{push\_back}.
\end{itemize}
Mixing them up is a common bug -- calling \texttt{reserve} when you
meant \texttt{resize} leaves the vector empty; calling
\texttt{resize} when you meant \texttt{reserve} creates $n$ junk
default-constructed elements that your \texttt{push\_back}s will
then \emph{append to}, giving you $2n$ total.
\end{warnbox}

\subsection{Related operations}

\begin{itemize}
    \item \texttt{v.clear()} -- empty the vector (size becomes 0;
          capacity usually untouched).
    \item \texttt{v.pop\_back()} -- remove the last element; shrinks
          size by 1.
    \item \texttt{v.empty()} -- returns \texttt{true} when size is
          0. Prefer \texttt{v.empty()} to \texttt{v.size() == 0}
          because on some containers (not \texttt{vector}, but good
          habit) \texttt{size()} is not $O(1)$.
\end{itemize}

\begin{warnbox}[\texttt{shrink\_to\_fit()} and the no-auto-shrink policy]
\texttt{std::vector} grows automatically but \emph{never shrinks
automatically.} If you push 10\,000 items then pop 9\,999, the vector's
size is~1 but its capacity is still $\geq 10\,000$ -- the unused buffer
sits there. \texttt{v.shrink\_to\_fit()} \emph{requests} a release, but
the standard explicitly allows the implementation to ignore it
(non-binding). To force the buffer down, swap with an empty vector:
\begin{lstlisting}
std::vector<int>().swap(v);   // capacity collapses to 0 (or near it)
\end{lstlisting}
This contrasts with Python's \texttt{list} (and some Java
\texttt{ArrayList} settings), which actively shrink the backing buffer
when too empty. The C++ default prioritizes reusing already-allocated
memory over reclaiming it -- usually what you want, occasionally not.
\end{warnbox}

\begin{ideabox}[The big picture for 1.6]
A \texttt{vector} has two knobs: \textbf{size} (what you can access)
and \textbf{capacity} (what has been allocated). \texttt{resize}
changes size; \texttt{reserve} changes only capacity;
\texttt{push\_back} changes size one at a time and lets capacity
grow geometrically behind the scenes. Get these three straight and
you will avoid most of the ``why is my vector empty / doubled /
reallocating constantly?'' bugs.
\end{ideabox}

\section{1.7 Vector push\_back, back, and pop\_back}

Section~1.6 already introduced \texttt{push\_back} as a way to grow a
vector one element at a time. The real value of this section is that
it completes the ``back of the vector'' interface -- \texttt{push\_back},
\texttt{back}, \texttt{pop\_back} -- which, taken together, turn a
plain \texttt{vector} into a perfectly good \textbf{stack}.

\subsection{The three back-end operations}

\begin{defnbox}[Back interface]
\begin{itemize}
    \item \texttt{v.push\_back(x)} -- append \texttt{x} to the end;
          size grows by 1.
    \item \texttt{v.back()} -- return a reference to the last element;
          size is unchanged.
    \item \texttt{v.pop\_back()} -- remove the last element; size
          shrinks by 1. Returns \texttt{void}.
    \item \texttt{v.front()} -- mirror of \texttt{back()} for the
          first element; also size-unchanged.
\end{itemize}
Every one of these is $O(1)$ (amortised, in \texttt{push\_back}'s
case).
\end{defnbox}

\begin{warnbox}[\texttt{pop\_back()} returns \texttt{void}]
A classic beginner mistake:
\begin{lstlisting}
std::cout << v.pop_back();   // ERROR: pop_back returns void
\end{lstlisting}
The textbook flags this but does not explain it, so here it is:
\texttt{pop\_back} separates ``give me the last value'' from ``remove
the last value.'' If you want both, call \texttt{back()} first, then
\texttt{pop\_back()}:
\begin{lstlisting}
int last = v.back();   // read the value
v.pop_back();          // then drop it
std::cout << last;
\end{lstlisting}
The reason the standard library splits them is \emph{exception
safety}: if returning the value could throw (copy-constructor of a
heavy object, out-of-memory), a ``pop-and-return'' would have to
somehow unwind a successful removal. Splitting the operations
sidesteps the problem -- and happens to be faster too, since
\texttt{back()} returns a reference without copying.
\end{warnbox}

\subsection{Vector as a stack}

\begin{ideabox}[Stack via vector]
\textbf{push, top, pop.} That is the entire interface of a
\textbf{stack} -- a last-in, first-out container. Rename the three
back operations accordingly:
\begin{itemize}
    \item \texttt{push\_back} $\equiv$ \texttt{push}
    \item \texttt{back}      $\equiv$ \texttt{top}
    \item \texttt{pop\_back} $\equiv$ \texttt{pop}
\end{itemize}
That is why a classic grocery-list example prints items in
\emph{reverse} order of entry -- ``Oranges, Apples, Bread, Juice''
comes out as ``Juice, Bread, Apples, Oranges.'' The last item pushed
is the first item popped.
\end{ideabox}

In fact the C++ standard library's \texttt{std::stack} is a thin
wrapper around exactly this pattern. You will meet it formally in
Chapter~4.

\subsection{Why the back is fast, the front is slow}

A common question: why is there no \texttt{push\_front} in the
headline interface?

\begin{ideabox}[Asymmetry of ends]
\texttt{push\_back} / \texttt{pop\_back} are $O(1)$ because they
touch only the last slot. Inserting or removing at the \emph{front}
would require shifting every other element one slot -- that is
$O(n)$. A \texttt{std::vector} actually does provide an
\texttt{insert(v.begin(), x)} operation, but it is expensive on
purpose. If you need fast front \emph{and} back access, use
\texttt{std::deque} (double-ended queue) or a linked list, which you
will meet in Chapter~4.
\end{ideabox}

This asymmetry is the reason \texttt{std::vector} makes a great stack
but a poor queue. A queue needs $O(1)$ on both ends; a vector does
not offer that.

\subsection{Worked example: building a list by reading input}

\begin{examplebox}[The CS 300 pattern]
Almost every project in this course builds a vector by reading an
unknown number of rows from a file. The shape is:
\begin{lstlisting}
std::vector<Bid> bids;
std::string line;
while (std::getline(file, line)) {
    bids.push_back(parseBid(line));
}
\end{lstlisting}
Because you do not know up front how many rows the CSV has, you
cannot \texttt{resize} in advance -- \texttt{push\_back} in a loop
is the natural fit. After the loop, \texttt{bids.size()} tells you
how many rows were read.
\end{examplebox}

\subsection{Small gotchas}

\begin{itemize}
    \item \textbf{\texttt{back()} on an empty vector} is undefined
          behaviour -- the reference it returns refers to nothing.
          Always test \texttt{!v.empty()} first when you are not
          sure.
    \item \textbf{\texttt{pop\_back()} on an empty vector} is also
          undefined behaviour.
    \item \textbf{A reference from \texttt{back()} is invalidated by
          a later \texttt{push\_back}} if the push triggers a
          reallocation. Read the value, use it, then push -- do not
          hold a reference across a push.
\end{itemize}

\begin{ideabox}[The big picture for 1.7]
The \texttt{push\_back} / \texttt{back} / \texttt{pop\_back} trio
gives you $O(1)$ access at the end of a vector and, together, form a
stack. Operations at the \emph{front} cost $O(n)$, so vectors are
LIFO-friendly and FIFO-hostile -- a useful distinction that shapes
which data structure you reach for in every later chapter.
\end{ideabox}

\section{1.8 Modifying, Copying, and Comparing Vectors}

Three distinct topics in one source-text section. The unifying idea,
and the thing to walk away with, is that a \texttt{std::vector} has
\textbf{value semantics}: assignment copies the data, equality compares
the data, and two vectors that were equal a moment ago can diverge
without affecting each other.

\subsection{Modifying elements in a loop}

You read an element with \texttt{v.at(i)}; you write one with
\texttt{v.at(i) = newValue}. There is nothing special about writing
during a loop -- the same index works for both directions:

\begin{lstlisting}
// Clamp every negative value to 0.
for (std::size_t i = 0; i < v.size(); ++i) {
    if (v.at(i) < 0) {
        v.at(i) = 0;
    }
}
\end{lstlisting}

\paragraph{The range-based form} is just as happy -- but the loop
variable has to be a \emph{reference} or you will be modifying a copy:

\begin{lstlisting}
for (int& x : v) {        // note the &
    if (x < 0) x = 0;
}
\end{lstlisting}

\begin{warnbox}[Forgetting the \texttt{\&} in a range-based for loop]
\texttt{for (int x : v)} copies each element into \texttt{x}. Writing
to \texttt{x} changes only the copy; \texttt{v} is untouched. Add
\texttt{\&} (\texttt{int\& x}) to modify in place, or use
\texttt{const int\& x} when you want to read without copying.
\end{warnbox}

\subsection{Shifting patterns -- the direction matters}

The two superficially similar left/right shifts in the textbook
exercises produce very different results because of the order the writes
happen:

\begin{examplebox}[Left shift -- overwrites ahead of the read]
\begin{lstlisting}
for (i = 0; i < 3; ++i) {
    v.at(i) = v.at(i + 1);
}
// {-55, -1, 0, 9}  ->  {-1, 0, 9, 9}
\end{lstlisting}
Each write happens \emph{before} its source is read, so the shift
works correctly from left to right. The last element is duplicated
because nothing is to its right to pull from.
\end{examplebox}

\begin{examplebox}[Right shift -- clobbers the data it is about to read]
\begin{lstlisting}
for (i = 0; i < 3; ++i) {
    v.at(i + 1) = v.at(i);
}
// {-55, -1, 0, 9}  ->  {-55, -55, -55, -55}
\end{lstlisting}
Here each iteration writes into the slot the \emph{next} iteration
needs to read. The original data gets smeared across the vector. To
right-shift correctly you have to walk from the end \emph{backwards}.
\end{examplebox}

\begin{ideabox}[The rule]
When modifying a vector in place using its own values, the iteration
direction determines whether you read ``before'' or ``after'' the
write. Walk in the direction that lets you read the original value
first, then write. If in doubt, copy to a scratch vector and read
from the copy.
\end{ideabox}

\subsection{Vector copy with \texttt{=}}

Assigning one vector to another copies every element:

\begin{lstlisting}
std::vector<int> origPrices = {10, 20, 30, 40};
std::vector<int> salePrices;

salePrices = origPrices;     // deep copy: salePrices is now {10, 20, 30, 40}
salePrices.at(2) = 27;       // does NOT affect origPrices
\end{lstlisting}

\begin{defnbox}[Value semantics]
A type has \textbf{value semantics} when copying a variable copies
the data, not a shared reference to it. After \texttt{b = a}, mutating
\texttt{b} has no effect on \texttt{a}. \texttt{std::vector},
\texttt{std::string}, and the built-in numeric types all work this way.
Contrast with Python or Java, where assigning one list/array variable
to another leaves both names pointing at the same object.
\end{defnbox}

\begin{ideabox}[Why this matters for performance]
Copying a vector of $n$ elements costs $O(n)$ time and allocates new
memory. That is invisible in a 4-element example but noticeable with
a million-row CSV. Two habits to build:
\begin{itemize}
    \item \textbf{Pass by reference} into functions:
          \texttt{void print(const std::vector<int>\& v)}. Without the
          \texttt{\&}, the whole vector is copied on every call.
    \item \textbf{Use \texttt{std::move}} when you are transferring
          ownership -- \texttt{b = std::move(a)} avoids the copy by
          stealing \texttt{a}'s internal buffer. \texttt{a} is left
          empty-but-valid.
\end{itemize}
\end{ideabox}

\subsection{Vector comparison with \texttt{==}}

\texttt{==} compares two vectors element-by-element: true if they
have the same size \emph{and} every element pair is equal.

\begin{lstlisting}
std::vector<int> x = {3, 4};
std::vector<int> y = {3, 4, 0, 7, 8};
std::vector<int> z = {3, 4, 0, 6, 8};

x == y;   // false: different sizes
y == z;   // false: same size, but index 3 differs
y == y;   // true
\end{lstlisting}

\begin{notebox}[\texttt{<}, \texttt{>}, \texttt{<=}, \texttt{>=} also work]
Not covered in the source text: the relational operators do
\textbf{lexicographic comparison}, the same way strings are ordered
in a dictionary. \texttt{\{1, 2, 9\} < \{1, 3, 0\}} because at the
first differing index the left side is smaller. Useful when you want
to sort a vector of vectors.
\end{notebox}

\begin{warnbox}[Never \texttt{==} floating-point vectors by accident]
Floating-point arithmetic rounds. \texttt{0.1 + 0.2 == 0.3} is
\texttt{false}. If you compare two \texttt{std::vector<double>}s that
came from any arithmetic at all, \texttt{==} will often surprise you.
Compare element-by-element with a tolerance (\texttt{abs(a - b) <
eps}) instead.
\end{warnbox}

\subsection{Related: does the vector's backing memory move?}

Worth knowing, since 1.6 introduced capacity:

\begin{itemize}
    \item \texttt{a = b} always \emph{at minimum} resets \texttt{a}'s
          size. If \texttt{b} is larger than \texttt{a}'s capacity, the
          internal buffer is reallocated.
    \item A \texttt{push\_back} that crosses the capacity threshold
          reallocates -- any references or iterators you were holding
          into the vector are now stale.
    \item Capacity \emph{grows} but does not shrink on its own; see
          the \texttt{shrink\_to\_fit} warnbox in \textsection 1.6.
\end{itemize}

\begin{ideabox}[The big picture for 1.8]
\texttt{std::vector} is a value type. \texttt{=} copies,
\texttt{==} compares, and modifications to one copy never leak into
the other. That makes reasoning about code easy -- but it also makes
unnecessary copies a real performance cost, so pass by
\texttt{const\&} whenever a function only needs to read a vector.
\end{ideabox}

\section{1.9 Swapping Two Variables}

Swapping looks trivial until you write it. The instinct
\texttt{x = y; y = x;} overwrites \texttt{x} before its original value is
ever used, so both variables end up holding \texttt{y}'s value and the old
\texttt{x} is gone. The fix is to save it somewhere first.

\begin{defnbox}[Swap]
To \textbf{swap} two variables is to assign each the other's value. If
\texttt{x = 33} and \texttt{y = 55}, after a swap \texttt{x = 55} and
\texttt{y = 33}. A correct swap requires a third location to temporarily
hold one value while the other is overwritten.
\end{defnbox}

\subsection{Why the naive version fails}

\begin{examplebox}[The broken swap]
\begin{verbatim}
int x = 33, y = 55;
x = y;   // x is now 55 -- the old 33 is GONE
y = x;   // y is now 55 (the new x), not 33
\end{verbatim}
After the first line there is no copy of \texttt{33} left anywhere in the
program. The second line cannot recover what was just overwritten.
\end{examplebox}

\subsection{The temp-variable pattern}

\begin{examplebox}[Correct swap with a temporary]
\begin{verbatim}
int x = 33, y = 55;
int temp = x;   // save x
x = y;          // clobber x -- fine, it's saved
y = temp;       // restore the saved value into y
\end{verbatim}
Three lines, one extra variable, always correct. The order matters:
save first, then do the overwrite, then use the saved copy.
\end{examplebox}

\subsection{C++ shortcut: \texttt{std::swap}}

C++ gives you the pattern for free in \texttt{<utility>} (or
\texttt{<algorithm>} in older standards), and you should almost always
use it instead of writing the temp-variable version by hand.

\begin{examplebox}[\texttt{std::swap}]
\begin{verbatim}
#include <utility>

int x = 33, y = 55;
std::swap(x, y);   // x == 55, y == 33
\end{verbatim}
It works on any type that supports move (or copy) assignment: \texttt{int},
\texttt{std::string}, \texttt{std::vector}, your own classes. For
\texttt{vector} and \texttt{string}, \texttt{std::swap} is O(1) -- it swaps
the internal pointers, not the elements.
\end{examplebox}

\begin{warnbox}[Do not name your variable \texttt{swap}]
If you write \texttt{using namespace std;} and then declare
\texttt{int swap;}, you shadow the function. Another reason to avoid
\texttt{using namespace std;} at global scope and to pick more specific
variable names.
\end{warnbox}

\subsection{Swapping list elements}

Most real swaps are between two positions in a vector, not between two
free-standing variables. Reversing a vector is the canonical example:
swap index \texttt{0} with \texttt{n-1}, index \texttt{1} with \texttt{n-2},
and so on, stopping at the middle.

\begin{examplebox}[Reverse a vector in place]
\begin{verbatim}
#include <vector>
#include <utility>

std::vector<int> v = {10, 20, 30, 40, 50};
for (size_t i = 0; i < v.size() / 2; ++i) {
    std::swap(v[i], v[v.size() - 1 - i]);
}
// v is now {50, 40, 30, 20, 10}
\end{verbatim}
The loop runs \texttt{n/2} times; integer division handles the odd-length
case correctly because the middle element swaps with itself and can be
skipped.
\end{examplebox}

\begin{warnbox}[Off-by-one: do not loop to \texttt{v.size()}]
If the loop bound is \texttt{v.size()} instead of \texttt{v.size() / 2},
every element gets swapped twice and the vector returns to its original
order. Reversing requires stopping at the midpoint.
\end{warnbox}

\begin{notebox}[\texttt{std::reverse} exists]
The standard library already has \texttt{std::reverse(v.begin(), v.end());}
in \texttt{<algorithm>}. Writing the loop yourself is useful for learning
the pattern, but production code should call the algorithm.
\end{notebox}

\begin{ideabox}[Big picture]
Swap is one of those building blocks that shows up inside almost every
sorting algorithm (selection sort, bubble sort, quicksort's partition
step, heapsort's sift-down). Understanding why the na\"ive two-line
version fails, and why a third storage slot is required, is the same
mental model you will reuse when reasoning about sort correctness later
in the course.
\end{ideabox}

\section{1.10 Debugging Example: Reversing a Vector}

The previous section gave the finished, correct reversal loop. This
section walks through the three bugs a first attempt typically hits, in
the order they surface, because the \emph{debugging process} is the real
lesson. Loop-plus-vector code is where most early C++ bugs live, and the
tracing technique used here works for everything later in the course.

\subsection{The buggy starting point}

Imagine writing this without thinking too hard:

\begin{examplebox}[Attempt 0 -- two bugs in three lines]
\begin{verbatim}
for (size_t i = 0; i < v.size(); ++i) {
    v.at(i) = v.at(v.size() - i);   // "swap"
}
\end{verbatim}
There are at least three things wrong here. Each one only becomes visible
after you fix the previous one, which is why tracing matters more than
guessing.
\end{examplebox}

\subsection{Bug 1: off-by-one on the mirror index}

\begin{warnbox}[\texttt{v.at(v.size() - i)} on the first iteration]
When \texttt{i == 0}, the expression is \texttt{v.at(v.size())}, which is
one past the end. \texttt{.at()} responds by throwing
\texttt{std::out\_of\_range} and the program aborts. The correct mirror
index is \texttt{v.size() - 1 - i}.
\end{warnbox}

\begin{notebox}[Why \texttt{.at()} is your friend here]
If the code had used \texttt{v[i]} the program would not have thrown; it
would have silently read whatever byte sat past the end of the buffer
and kept going. That is undefined behavior. \texttt{.at()} turns a silent
memory bug into a loud, diagnosable exception. Use it while debugging.
\end{notebox}

\subsection{Bug 2: assignment is not a swap}

After fixing the index, the loop runs without crashing but produces the
wrong output -- values appear duplicated and others disappear:

\begin{examplebox}[Tracing attempt 1]
Input \texttt{10 20 30 40 50 60 70 80}, loop body
\texttt{v.at(i) = v.at(v.size() - 1 - i);}
\begin{verbatim}
i=0: v[0] = v[7]  ->  80 20 30 40 50 60 70 80
i=1: v[1] = v[6]  ->  80 70 30 40 50 60 70 80
i=2: v[2] = v[5]  ->  80 70 60 40 50 60 70 80
i=3: v[3] = v[4]  ->  80 70 60 50 50 60 70 80   // 40 is gone
i=4: v[4] = v[3]  ->  80 70 60 50 50 60 70 80
\end{verbatim}
The original value at \texttt{v[0]} was overwritten before it was
copied anywhere. This is the same bug as the na\"ive two-line swap from
\textsection 1.9: assignment copies one way only.
\end{examplebox}

The fix is the temp-variable swap (or \texttt{std::swap}):

\begin{verbatim}
int tmp = v.at(i);
v.at(i) = v.at(v.size() - 1 - i);
v.at(v.size() - 1 - i) = tmp;
\end{verbatim}

\subsection{Bug 3: the loop runs twice as long as it should}

After fixing the swap, the output comes out identical to the input. The
swap is correct this time; the loop just runs too many iterations:

\begin{examplebox}[Tracing attempt 2]
\begin{verbatim}
i=0: swap v[0] and v[7]  ->  80 20 30 40 50 60 70 10
i=1: swap v[1] and v[6]  ->  80 70 30 40 50 60 20 10
i=2: swap v[2] and v[5]  ->  80 70 60 40 50 30 20 10
i=3: swap v[3] and v[4]  ->  80 70 60 50 40 30 20 10  // reversed
i=4: swap v[4] and v[3]  ->  80 70 60 40 50 30 20 10  // un-reversed
i=5: swap v[5] and v[2]  ->  ...                      // etc.
\end{verbatim}
Once the midpoint is crossed, every further swap un-does an earlier one.
The loop bound must be \texttt{i < v.size() / 2}.
\end{examplebox}

\begin{warnbox}[Integer division quietly handles the odd-length case]
If \texttt{v.size()} is 5, \texttt{v.size() / 2} is 2. The loop swaps
\texttt{(0, 4)} and \texttt{(1, 3)} and stops; the middle element at
index 2 swaps with itself and is correctly left alone. No extra branch
needed.
\end{warnbox}

\subsection{The debugging methodology, not the answer}

The final correct version was already given in \textsection 1.9. What
matters here is the \emph{technique}:

\begin{enumerate}[leftmargin=*]
    \item \textbf{Read the error.} \texttt{out\_of\_range} pointed straight
          at an index expression. Don't skip the stack trace.
    \item \textbf{Trace by hand.} Write out the vector contents after each
          iteration on paper or in a comment. Bugs 2 and 3 are invisible
          from the source but obvious after three lines of trace output.
    \item \textbf{Fix one bug at a time.} Each fix exposes the next bug.
          Changing three things at once means you will not know which
          change did what.
    \item \textbf{Use \texttt{.at()} during development.} Turn silent
          out-of-bounds into exceptions. Switch to \texttt{[]} later if
          you have measured that bounds checking is a bottleneck (it
          rarely is).
\end{enumerate}

\begin{ideabox}[Big picture]
Reversing a vector is not the skill being taught; the skill is the
tight loop of \emph{predict, run, observe, diff}. When a loop does the
wrong thing, resist the urge to stare at the code. Instrument it --
print the state after each iteration, or step through in a debugger --
and compare what you see to what you expected. That habit scales from
three-line swap loops all the way up to the sorts, searches, and graph
algorithms later in the course.
\end{ideabox}

\section{1.11 Arrays vs. Vectors}

C++ inherits the raw \textbf{C array} and also provides
\texttt{std::vector}. They look similar on the surface -- indexed access,
contiguous storage -- but they have very different safety and feature
profiles. For this course the rule is simple: \textbf{use
\texttt{std::vector} unless a concrete reason forces otherwise}.

\subsection{Syntax side-by-side}

\begin{examplebox}[Declaration and access]
\begin{verbatim}
// C-style array: fixed size at compile time
int a[5];                 // 5 ints, uninitialized!
a[0] = 42;
int x = a[0];

// std::vector: size known at runtime, grows if needed
#include <vector>
std::vector<int> v(5);    // 5 ints, all zero-initialized
v.at(0) = 42;             // bounds-checked
int y = v.at(0);
int z = v[0];             // also works, NOT bounds-checked
\end{verbatim}
\end{examplebox}

\subsection{Why arrays are dangerous}

An out-of-range write to a C array does not throw, does not warn, does
not even necessarily crash. It just scribbles over whatever happens to
live at that memory address.

\begin{warnbox}[The classic memory-corruption bug]
\begin{verbatim}
int weights[3];
int age = 44;

weights[0] = 122;
weights[1] = 119;
weights[2] = 117;
weights[3] = 199;   // out of range -- writes into age!

std::cout << age;   // prints 199 (on many compilers)
\end{verbatim}
The compiler lays \texttt{weights} and \texttt{age} next to each other
in memory. Writing \texttt{weights[3]} lands on \texttt{age}'s bytes.
Nothing stops it. The program appears correct until something later
reads \texttt{age} and gets a nonsense value -- possibly hours of
debugging away from the real bug.
\end{warnbox}

With \texttt{std::vector} the same mistake using \texttt{.at()} throws
\texttt{std::out\_of\_range} immediately, pointing at the exact line.

\subsection{Feature comparison}

\begin{center}
\begin{tabular}{lll}
\textbf{Feature} & \textbf{C array} & \texttt{std::vector} \\
\hline
Size fixed at & compile time\textsuperscript{*} & runtime \\
Can grow / shrink & no & yes \\
Knows its own size & no & yes (\texttt{.size()}) \\
Bounds checking available & no & yes (\texttt{.at()}) \\
Copy with \texttt{=} & no (decays to pointer) & yes (deep copy) \\
Pass to function & as pointer + length & as value or \texttt{const\&} \\
Compare with \texttt{==} & no (pointer compare) & yes (element-wise) \\
\end{tabular}
\end{center}

\noindent\small\textsuperscript{*}C99 variable-length arrays exist but
are a GCC extension in C++ and should be avoided.\normalsize

\begin{defnbox}[Array decay]
When a C array is passed to a function or assigned to another variable,
it silently converts to a bare pointer to its first element. The size
information is \textbf{lost}. This is why C-style APIs always pass a
separate length parameter (\texttt{void f(int* a, size\_t n)}), and why
\texttt{sizeof(a) / sizeof(a[0])} only works in the scope where the
array was declared.
\end{defnbox}

\subsection{When arrays still show up}

You will not eliminate arrays from your C++ code -- they appear in
places you cannot avoid:

\begin{itemize}[leftmargin=*]
    \item \textbf{String literals} are arrays: \texttt{"hello"} is a
          \texttt{const char[6]}.
    \item \texttt{argv} in \texttt{main(int argc, char* argv[])} is a
          C-style array of C-style strings.
    \item \textbf{Legacy C APIs} (POSIX, Win32, most game engines) take
          pointer-plus-length pairs.
    \item \textbf{Embedded / performance} code sometimes needs the
          stack-only, zero-overhead guarantee of a raw array.
\end{itemize}

For the \textbf{fixed-size} cases where you would be tempted to use a
C array, prefer \texttt{std::array<T, N>} -- it has a compile-time
size like a C array but carries \texttt{.size()}, \texttt{.at()},
iterators, and value semantics like a vector.

\begin{examplebox}[\texttt{std::array} as the safe fixed-size choice]
\begin{verbatim}
#include <array>

std::array<int, 3> weights = {122, 119, 117};
weights.at(3) = 199;   // throws std::out_of_range
int n = weights.size();  // 3, known at compile time
\end{verbatim}
\end{examplebox}

\begin{warnbox}[\texttt{arrayList.size()} does not compile]
A common reflex after using \texttt{vector} is to write
\texttt{a.size()} on a C array. Arrays are not objects; they have no
member functions. The language-level size of a declared array is
\texttt{sizeof(a) / sizeof(a[0])}, and even that is wrong after decay.
If you want \texttt{.size()}, you want \texttt{std::array} or
\texttt{std::vector}.
\end{warnbox}

\begin{ideabox}[Big picture]
Arrays are what the hardware actually provides: a contiguous block of
memory, indexed by offset, with no bookkeeping. Vectors are what you
usually want: the same contiguous block, but with a size, growth,
bounds checking, and value semantics layered on top. The extra
bookkeeping costs essentially nothing, and the safety is worth it in
every course project. Default to \texttt{std::vector}; reach for
\texttt{std::array} when the size is fixed; reach for a raw array only
when an external API forces it.
\end{ideabox}

\section{1.12 Two-Dimensional Arrays}

A 2D array is a table: \texttt{R} rows $\times$ \texttt{C} columns, for
\texttt{R * C} elements total. The language lets you write
\texttt{a[r][c]} and pretend you are indexing a grid, but underneath,
memory is flat. Knowing how the rows are laid out explains cache
performance, how to pass 2D data to functions, and why
\texttt{std::vector<std::vector<T>>} is not actually a 2D array.

\begin{defnbox}[Row-major order]
C and C++ store 2D arrays in \textbf{row-major} order: all of row 0
first, then all of row 1, then row 2, and so on, in one contiguous
block of memory. Given \texttt{int a[R][C]}, element \texttt{a[r][c]}
lives at offset \texttt{r * C + c} from the start of the array. The
\emph{rightmost} index changes fastest as you walk memory.
\end{defnbox}

\subsection{Declaring and using}

\begin{examplebox}[Basic 2D array]
\begin{verbatim}
int grid[2][3];              // 2 rows, 3 cols, 6 ints total

grid[0][0] = 55;
grid[1][2] = 99;

// Brace-initialization makes the shape visible:
int numVals[2][3] = {
    {22, 44, 66},   // row 0
    {97, 98, 99}    // row 1
};
\end{verbatim}
\end{examplebox}

\subsection{Iterating: the loop order matters}

The idiomatic nested loop walks the outer (row) index first, the inner
(column) index second:

\begin{examplebox}[Standard nested loop]
\begin{verbatim}
for (int r = 0; r < R; ++r) {
    for (int c = 0; c < C; ++c) {
        std::cout << grid[r][c] << ' ';
    }
    std::cout << '\n';
}
\end{verbatim}
\end{examplebox}

\begin{notebox}[Why row-first is faster than column-first]
Walking \texttt{r} on the outside and \texttt{c} on the inside visits
memory \emph{sequentially}: \texttt{grid[0][0], grid[0][1], grid[0][2],
grid[1][0], \ldots} The CPU prefetches the next cache line before you
ask for it. Reversing the loops (\texttt{c} outside, \texttt{r} inside)
jumps by \texttt{C} elements each access, thrashing the cache. For a
small 3$\times$3 it doesn't matter; for a 10\,000$\times$10\,000 matrix
the two orders can differ by a factor of ten.
\end{notebox}

\subsection{Passing a 2D array to a function}

This is where raw 2D arrays bite. The compiler needs to know the row
length to compute \texttt{r * C + c}, so all dimensions \emph{except the
first} must be spelled out in the parameter type:

\begin{examplebox}[Three equivalent signatures]
\begin{verbatim}
void f(int a[2][3], int rows);   // literal dimensions
void f(int a[][3], int rows);    // first dim can be omitted
void f(int (*a)[3], int rows);   // pointer to row (same thing)
\end{verbatim}
\end{examplebox}

\begin{warnbox}[\texttt{void f(int** a)} does NOT take a 2D array]
A 2D array is \emph{not} an array of pointers. \texttt{int a[3][4]} is
12 contiguous ints, not 3 pointers to 4 ints. Passing it where
\texttt{int**} is expected compiles only after a reinterpret cast and
will then read garbage. If you need a function that accepts any row
length, use templates or pass a pointer plus both dimensions.
\end{warnbox}

\subsection{The C++ alternative: \texttt{std::vector} of \texttt{vector}}

For runtime-sized 2D data the common idiom is a vector of vectors.

\begin{examplebox}[2D as a vector-of-vectors]
\begin{verbatim}
#include <vector>

int R = 4, C = 5;
std::vector<std::vector<int>> grid(R, std::vector<int>(C, 0));

grid[1][2] = 99;          // same syntax as a 2D array
grid.size();              // R
grid[0].size();           // C
\end{verbatim}
This is flexible -- rows can even have different lengths -- but the
rows are \emph{not} contiguous in memory. Each inner vector is a
separate heap allocation, and cache behavior suffers accordingly.
\end{examplebox}

\begin{ideabox}[Flat 1D storage is often better]
For true matrix work the trick is to store one vector of
\texttt{R * C} elements and compute the index yourself:
\begin{verbatim}
std::vector<int> grid(R * C);
auto at = [&](int r, int c) -> int& { return grid[r * C + c]; };
at(1, 2) = 99;
\end{verbatim}
One allocation, contiguous memory, same big-O as a raw 2D array, plus
\texttt{.size()} and bounds-checked iteration. This is how most
numerical C++ code represents 2D data.
\end{ideabox}

\subsection{Higher dimensions}

\texttt{int a[2][3][5]} declares 30 elements; \texttt{int a[100][100][5][3]}
declares 150\,000. Dimensions multiply, and stack-allocated arrays can
blow the stack surprisingly fast. If any dimension is large, allocate
with \texttt{std::vector} (heap) rather than a raw array.

\begin{ideabox}[Big picture]
A 2D array is a convenient fiction the compiler maintains over a flat
memory block. Understanding row-major order tells you which loop order
is fast, why the column dimensions must be in the function signature,
and why \texttt{vector<vector<T>>} and a \texttt{vector<T>} with manual
indexing have different performance characteristics despite looking
similar. Most course work is small enough that any of these options
works; the habits you build here pay off when the matrices grow.
\end{ideabox}

\section{1.13 Char Arrays / C Strings}

Before \texttt{std::string} existed, a \textbf{string} in C was just an
array of \texttt{char}s ending with a special sentinel byte. C++ still
supports this representation -- it's what you get from string literals,
command-line arguments, and every legacy C API you will ever link
against -- so it is worth understanding even though modern C++ code
should use \texttt{std::string}.

\begin{defnbox}[Null-terminated string (``C string'')]
A \textbf{C string} is a sequence of \texttt{char} values stored in a
\texttt{char} array, terminated by the \textbf{null character}
\texttt{'\textbackslash0'} (ASCII code 0). The length of the string is
not stored anywhere -- it is \emph{determined} by scanning for the null
byte. A \texttt{char[20]} array can therefore hold strings of length 0
through 19: one slot is always reserved for the terminator.
\end{defnbox}

\subsection{Declaring a C string}

\begin{examplebox}[Three equivalent declarations]
\begin{verbatim}
char movie[20] = "Star Wars";  // explicit size, rest zero-filled
char name[] = "Hellen";        // size inferred: 7 ('H','e','l','l','e','n','\0')
char city[5] = {'A','l','a','n','\0'};  // character-by-character
\end{verbatim}
In every case the compiler (or you) must make room for the trailing
\texttt{'\textbackslash0'}. For the literal \texttt{"Amy"},
\texttt{char name[4]} is the minimum -- three letters plus the null.
\end{examplebox}

\begin{warnbox}[\texttt{char name[3] = "Amy";} silently drops the null]
Some older compilers accept this and just don't write a terminator;
newer ones will warn or error. Every subsequent \texttt{cout << name;}
or \texttt{strlen(name)} then reads past the end of the array.
\textbf{Always size the array for the literal \emph{plus one}.}
\end{warnbox}

\subsection{Why printing ``just works''}

\texttt{std::cout << some\_c\_string} does not print 20 characters. It
prints characters one at a time until it sees \texttt{'\textbackslash0'}.
That convention runs through every C string API -- \texttt{printf},
\texttt{strlen}, \texttt{strcpy}, file I/O -- which is why a missing or
overwritten null byte turns into a memory-corruption bug instead of a
tidy error.

\subsection{Three common bugs}

\begin{warnbox}[1. Looping to the array size instead of to the null]
\begin{verbatim}
char s[20] = "hello";
for (int i = 0; i < 20; ++i)  cout << s[i];   // WRONG
\end{verbatim}
Prints \texttt{hello} then 14 bytes of whatever happens to sit after
the visible text. Correct loop:
\begin{verbatim}
for (int i = 0; s[i] != '\0'; ++i)  cout << s[i];
\end{verbatim}
\end{warnbox}

\begin{warnbox}[2. Looping past the array entirely]
\begin{verbatim}
for (int i = 0; i < 30; ++i)  cout << s[i];   // UB
\end{verbatim}
Reads memory that is not part of \texttt{s} at all. May print garbage,
may crash, may appear to work. This is undefined behavior.
\end{warnbox}

\begin{warnbox}[3. Overwriting the null terminator]
\begin{verbatim}
char name[5] = "Alan";   // {'A','l','a','n','\0'}
name[4] = '!';           // null is gone
cout << name;            // prints "Alan!" then keeps going...
\end{verbatim}
The next \texttt{cout} walks off the end of the array looking for a
\texttt{'\textbackslash0'}. It might hit one quickly (harmless-looking
output) or run for a long time printing junk (obvious garbage) -- the
bug is nondeterministic.
\end{warnbox}

\subsection{You cannot assign to a C string}

\begin{examplebox}[The assignment that doesn't assign]
\begin{verbatim}
char title[20] = "Star Wars";
title = "Indiana Jones";   // ERROR: array name is not an lvalue
\end{verbatim}
A C array's name is not a variable you can reassign. To \emph{copy}
characters into an existing buffer, use \texttt{strcpy} or
\texttt{strncpy} from \texttt{<cstring>}:
\begin{verbatim}
#include <cstring>
std::strcpy(title, "Indiana Jones");  // title must be big enough!
\end{verbatim}
\end{examplebox}

\begin{warnbox}[\texttt{strcpy} has no idea how big the destination is]
Copying a 30-character string into a \texttt{char[20]} compiles
cleanly, writes 10 bytes past the end of the array, and silently
corrupts memory -- classic buffer overflow, historically the root of
countless security vulnerabilities. Use \texttt{strncpy}, enforce your
own bounds check, or (better) use \texttt{std::string}.
\end{warnbox}

\subsection{C strings vs. \texttt{std::string}}

\begin{center}
\begin{tabular}{lll}
& \textbf{C string} (\texttt{char[]}) & \texttt{std::string} \\
\hline
Header & \texttt{<cstring>} & \texttt{<string>} \\
Stores length & no (scan for \texttt{\textbackslash0}) & yes (\texttt{.size()}) \\
Assignable with \texttt{=} & no & yes \\
Concat with \texttt{+} & no & yes \\
Resizes automatically & no & yes \\
Bounds-checked access & no & \texttt{.at()} \\
\end{tabular}
\end{center}

\begin{notebox}[\texttt{std::string} gives you a C string on demand]
Every C++ \texttt{string} carries a null-terminated char buffer
internally. Call \texttt{s.c\_str()} to get a \texttt{const char*}
suitable for passing to a legacy C API:
\begin{verbatim}
std::string path = "data.txt";
std::FILE* f = std::fopen(path.c_str(), "r");
\end{verbatim}
This is how you bridge modern C++ code to C libraries.
\end{notebox}

\begin{ideabox}[Big picture]
The null terminator is one of the oldest and leakiest abstractions in
the C language -- fast and compact, but every single string operation
has to trust that the terminator is where the caller said it would be.
Every buffer overflow in the classic sense comes from that trust being
broken. \texttt{std::string} wraps the same raw buffer with an explicit
size, safer copy/assign, and automatic resizing, which is why C++ code
uses it almost everywhere. You learn C strings because they are still
out there -- in \texttt{argv}, in POSIX APIs, in file paths -- not
because you should write new code that uses them.
\end{ideabox}

\section{1.14 String Library Functions}

Because C strings do not know their own length and do not support
\texttt{=} or \texttt{==}, the standard library ships a family of free
functions that take raw \texttt{char*} pointers and walk them looking
for the null terminator. They all live in \texttt{<cstring>} (the C++
name) or \texttt{<string.h>} (the C name). You will recognize most of
them from any C tutorial.

\begin{defnbox}[The naming convention]
Functions starting with \texttt{str} work on null-terminated C strings.
Functions starting with \texttt{mem} work on raw byte buffers with an
explicit length. Functions with an \texttt{n} in the middle
(\texttt{strncpy}, \texttt{strncat}, \texttt{strncmp}) take an extra
length argument and are the safer variant of their bare cousin.
\end{defnbox}

\subsection{Modifying strings}

\begin{center}
\small
\begin{tabular}{ll}
\texttt{strcpy(dst, src)} & Copy \texttt{src} (through null) into \texttt{dst}. \\
\texttt{strncpy(dst, src, n)} & Copy up to \texttt{n} chars. \\
\texttt{strcat(dst, src)} & Append \texttt{src} to the end of \texttt{dst}. \\
\texttt{strncat(dst, src, n)} & Append up to \texttt{n} chars, then a null. \\
\end{tabular}
\normalsize
\end{center}

\begin{examplebox}[\texttt{strcpy} / \texttt{strcat} basics]
\begin{verbatim}
#include <cstring>
char org[100] = "United Nations";
char user[20] = "UNICEF";
char target[10];

std::strcpy(target, user);   // target == "UNICEF"
std::strcat(org, user);      // org    == "United NationsUNICEF"
\end{verbatim}
Destination comes first, source second -- same order as
\texttt{target = source;} in regular C++. Swapping the arguments is a
classic beginner bug and the compiler cannot catch it.
\end{examplebox}

\begin{warnbox}[\texttt{strcpy} does not check the destination size]
\begin{verbatim}
char target[10];
std::strcpy(target, org);   // org is 14 chars -- writes 4 past the end!
\end{verbatim}
No exception, no error -- just silent memory corruption. This is the
canonical buffer-overflow bug. Use \texttt{strncpy(dst, src, sizeof
dst)} if you must, and note even \texttt{strncpy} does not guarantee
a trailing null when the source is exactly \texttt{n} characters long.
\end{warnbox}

\begin{warnbox}[Array-name assignment still does not copy]
\begin{verbatim}
target = user;     // ERROR: array name is not an lvalue
\end{verbatim}
The \emph{only} place \texttt{=} copies a C string is in the
declaration itself, because the compiler is initializing the array,
not performing a real assignment:
\begin{verbatim}
char user[20] = "UNICEF";    // OK, compile-time initialization
\end{verbatim}
\end{warnbox}

\subsection{Inspecting strings}

\begin{center}
\small
\begin{tabular}{ll}
\texttt{strlen(s)} & Returns number of chars before the null (type \texttt{size\_t}). \\
\texttt{strcmp(a, b)} & $<0$ if \texttt{a < b}, 0 if equal, $>0$ if \texttt{a > b}. \\
\texttt{strncmp(a, b, n)} & Same, limited to first \texttt{n} chars. \\
\texttt{strchr(s, ch)} & Pointer to first \texttt{ch} in \texttt{s}, or \texttt{NULL}. \\
\texttt{strstr(hay, needle)} & Pointer to first occurrence of \texttt{needle}, or \texttt{NULL}. \\
\end{tabular}
\normalsize
\end{center}

\begin{examplebox}[Length, search, compare]
\begin{verbatim}
size_t n = std::strlen("Alan Turing");     // 11

if (std::strchr(name, ' ') != nullptr) { /* has a space */ }

if (std::strcmp(s1, s2) == 0)  { /* equal */ }
if (std::strcmp(s1, s2) < 0)   { /* s1 comes first */ }
\end{verbatim}
\end{examplebox}

\subsection{The three \texttt{strcmp} traps}

These come up constantly and none of them are syntax errors -- the
compiler will not catch any of them.

\begin{warnbox}[1. Using \texttt{==} on C strings]
\begin{verbatim}
if (s1 == s2) { ... }   // compares POINTER addresses, not content
\end{verbatim}
Two character arrays will almost never share an address, so this
branch is effectively dead. Silent logic bug.
\end{warnbox}

\begin{warnbox}[2. Forgetting \texttt{== 0}]
\begin{verbatim}
if (std::strcmp(s1, s2)) { ... }   // backwards!
\end{verbatim}
\texttt{strcmp} returns \textbf{0 when equal}. Treating the result as a
boolean reads ``truthy means equal,'' which is the opposite of what it
actually means. Always write \texttt{== 0} explicitly -- your future
self, reading the code cold, should not have to remember the
convention.
\end{warnbox}

\begin{warnbox}[3. \texttt{strlen} on a buffer that has no null]
\begin{verbatim}
char buf[10];             // no initializer
size_t n = std::strlen(buf);   // walks memory until it hits any 0 byte
\end{verbatim}
Undefined behavior. If you declare a C-string buffer, either
initialize it (\texttt{char buf[10] = "";}) or do not call
string-library functions on it until you have written real content
plus a terminator.
\end{warnbox}

\subsection{A common pattern: transform characters in place}

\begin{examplebox}[Replace spaces with underscores]
\begin{verbatim}
#include <cstring>
char name[15] = "Alan Turing";

for (size_t i = 0; i < std::strlen(name); ++i) {
    if (name[i] == ' ') name[i] = '_';
}
\end{verbatim}
This loop \textbf{re-scans the whole string every iteration} because
\texttt{strlen} is O(n). For long strings, hoist it:
\begin{verbatim}
size_t len = std::strlen(name);
for (size_t i = 0; i < len; ++i) { ... }
\end{verbatim}
Or stop at the null byte directly, which is the idiomatic form:
\begin{verbatim}
for (size_t i = 0; name[i] != '\0'; ++i) { ... }
\end{verbatim}
\end{examplebox}

\subsection{The modern C++ translation}

Every function in this section has a cleaner \texttt{std::string}
equivalent that does not buffer-overflow and knows its own size.

\begin{center}
\small
\begin{tabular}{ll}
\textbf{C} & \textbf{C++} (\texttt{std::string}) \\
\hline
\texttt{strcpy(dst, src)} & \texttt{dst = src;} \\
\texttt{strcat(dst, src)} & \texttt{dst += src;} \\
\texttt{strlen(s)} & \texttt{s.size()} \\
\texttt{strcmp(a, b) == 0} & \texttt{a == b} \\
\texttt{strcmp(a, b) < 0} & \texttt{a < b} \\
\texttt{strchr(s, ch)} & \texttt{s.find(ch)} \\
\texttt{strstr(h, n)} & \texttt{h.find(n)} \\
\end{tabular}
\normalsize
\end{center}

\begin{ideabox}[Big picture]
The \texttt{cstring} library is doing, function by function, what
\texttt{std::string} does as member operations. Every bug that
\texttt{strcpy}, \texttt{strcmp}, and \texttt{strlen} are famous for
(overflow, wrong-direction comparison, missing null) disappears when
you let the string type track its own size. Learn these functions so
you can read legacy code and glue against C APIs -- then write new C++
code with \texttt{std::string}.
\end{ideabox}

\section{1.15 Char Library Functions: \texttt{cctype}}

\texttt{<cctype>} is the C++ header for classifying and converting
individual \texttt{char} values -- ``is this letter a digit?'' ``convert
this to uppercase.'' It's tiny, and every function takes an
\texttt{int}, returns an \texttt{int}, and is named
\texttt{is\textit{something}} or \texttt{to\textit{something}}. Learn
the short list and you have 95\% of the character-manipulation toolbox.

\begin{defnbox}[Why it's \texttt{cctype} not \texttt{ctype}]
C headers in C++ gain a leading \texttt{c} and lose the \texttt{.h}:
\texttt{<ctype.h>} becomes \texttt{<cctype>},
\texttt{<string.h>} becomes \texttt{<cstring>},
\texttt{<stdio.h>} becomes \texttt{<cstdio>}. The C++ version puts
everything in namespace \texttt{std} (so prefer \texttt{std::isalpha}
over bare \texttt{isalpha}).
\end{defnbox}

\subsection{Classification: ``is this character a \ldots?''}

All of these return an \texttt{int} -- nonzero for true, 0 for false.
You will usually use them in an \texttt{if} condition directly.

\begin{center}
\small
\begin{tabular}{ll}
\texttt{isalpha(c)} & alphabetic: \texttt{a}--\texttt{z} or \texttt{A}--\texttt{Z} \\
\texttt{isdigit(c)} & decimal digit: \texttt{0}--\texttt{9} \\
\texttt{isalnum(c)} & \texttt{isalpha} or \texttt{isdigit} \\
\texttt{isspace(c)} & space, tab, newline, CR, vertical tab, form-feed \\
\texttt{isblank(c)} & just space or tab \\
\texttt{islower(c)} / \texttt{isupper(c)} & \texttt{a}--\texttt{z} / \texttt{A}--\texttt{Z} \\
\texttt{isxdigit(c)} & hex digit: \texttt{0}--\texttt{9}, \texttt{a}--\texttt{f}, \texttt{A}--\texttt{F} \\
\texttt{ispunct(c)} & printable, not space, not alphanumeric \\
\texttt{isprint(c)} & printable (includes space) \\
\texttt{iscntrl(c)} & control char -- the complement of \texttt{isprint} \\
\end{tabular}
\normalsize
\end{center}

\begin{examplebox}[Classifying characters in a string]
\begin{verbatim}
#include <cctype>
char s[] = "Hey9! Go";

std::isalpha(s[0]);  // 'H' -> nonzero (true)
std::isdigit(s[3]);  // '9' -> nonzero (true)
std::isspace(s[5]);  // ' ' -> nonzero (true)
std::ispunct(s[4]);  // '!' -> nonzero (true)
\end{verbatim}
\end{examplebox}

\subsection{Conversion: \texttt{toupper} and \texttt{tolower}}

Both take a character and return a (possibly different) character. If
the input isn't a letter of the relevant case, the function returns it
unchanged -- so you can apply it to every character in a string
without branching.

\begin{examplebox}[Uppercase the whole string in place]
\begin{verbatim}
#include <cctype>
char s[] = "Hey9! Go";

for (size_t i = 0; s[i] != '\0'; ++i) {
    s[i] = std::toupper(static_cast<unsigned char>(s[i]));
}
// s is now "HEY9! GO"
\end{verbatim}
\end{examplebox}

\begin{warnbox}[The \texttt{unsigned char} cast is not paranoia]
\texttt{isalpha}, \texttt{toupper}, and the rest are specified to take
an \texttt{int} whose value is either \texttt{EOF} (\texttt{-1}) or a
representation of an \texttt{unsigned char}. On platforms where
\texttt{char} is \emph{signed} (most of them), passing a character
with the high bit set -- anything above ASCII 127, e.g. accented
letters -- produces a negative \texttt{int} and undefined behavior.
Always cast: \texttt{std::toupper(static\_cast<unsigned char>(c))}.
This one line prevents a nasty class of bugs you will only see on
non-English input.
\end{warnbox}

\subsection{A concrete use case: case-insensitive compare}

\texttt{strcmp} is case-sensitive. \texttt{"Apple"} and
\texttt{"apple"} are not equal to it. The standard trick is to
lowercase both sides and compare.

\begin{examplebox}[Case-insensitive equality for C strings]
\begin{verbatim}
bool equal_ci(const char* a, const char* b) {
    while (*a && *b) {
        char ca = std::tolower(static_cast<unsigned char>(*a));
        char cb = std::tolower(static_cast<unsigned char>(*b));
        if (ca != cb) return false;
        ++a; ++b;
    }
    return *a == *b;   // both must now be '\0'
}
\end{verbatim}
\end{examplebox}

\subsection{The C++20 alternative: ranges + lambdas}

When you are already using \texttt{std::string} (and you should be),
the loop gets shorter and the cast disappears into the algorithm.

\begin{examplebox}[Uppercase with \texttt{std::transform}]
\begin{verbatim}
#include <algorithm>
#include <cctype>
#include <string>

std::string s = "Hey9! Go";
std::transform(s.begin(), s.end(), s.begin(),
    [](unsigned char c){ return std::toupper(c); });
// s == "HEY9! GO"
\end{verbatim}
The lambda parameter is declared \texttt{unsigned char} so the cast
happens automatically at the call site -- cleaner than scattering
\texttt{static\_cast} through a loop body.
\end{examplebox}

\subsection{Parsing pattern: extract digits from a string}

This pops up in every intro assignment that asks you to pull a phone
number, zip code, or ID out of user input:

\begin{examplebox}[Keep only digits]
\begin{verbatim}
#include <cctype>
#include <string>

std::string phone = "(603) 555-1212";
std::string digits;
for (unsigned char c : phone) {
    if (std::isdigit(c)) digits.push_back(c);
}
// digits == "6035551212"
\end{verbatim}
\end{examplebox}

\begin{ideabox}[Big picture]
\texttt{cctype} is the character-level counterpart to \texttt{cstring}:
one handles single chars, the other walks null-terminated arrays of
them. Together they cover classic C text processing. Modern C++ adds
\texttt{std::string}, algorithms, and regular expressions on top, but
the \texttt{is*} / \texttt{to*} functions still show up inside those
higher-level tools -- so learning the short list here pays off
immediately, and does not become obsolete when you graduate to
\texttt{std::string}.
\end{ideabox}

\section*{Summary}

Chapter 1 is a whirlwind review of the C++ building blocks every later
data-structures topic assumes. The threads worth carrying forward:

\begin{itemize}[leftmargin=*]
    \item \textbf{Prefer safe abstractions.} \texttt{std::vector} over
    raw arrays, \texttt{std::string} over \texttt{char[]}, \texttt{.at()}
    over \texttt{[]} while debugging. Every C-flavored alternative has
    a silent-failure mode that costs more to debug than the
    bounds-checked version costs to run.
    \item \textbf{Indices are contracts.} \texttt{size\_t} is unsigned;
    \texttt{v.size() - 1} underflows catastrophically on an empty
    vector; a 2D index is \texttt{r * C + c} under the hood. Treat
    every index expression as something to verify, not guess.
    \item \textbf{Copy is not free.} Vector copy is O(n); passing a big
    container by value does a deep copy. Use \texttt{const\&} when you
    only read, \texttt{\&} when you mutate, and move semantics when
    ownership transfers. Value semantics are the reason all of this
    feels sane compared to C.
    \item \textbf{Algorithms over loops.} \texttt{accumulate},
    \texttt{max\_element}, \texttt{count\_if}, \texttt{find},
    \texttt{reverse}, \texttt{swap}, \texttt{transform} -- the
    \texttt{<algorithm>} header already has the loop you were about to
    write. Reach for it first.
    \item \textbf{Debug by tracing, not by staring.} When a loop
    produces the wrong output, print the state after every iteration.
    Three bugs in a ten-line reversal loop were only findable that way.
    \item \textbf{Vectors are the runway for the whole course.} Stacks,
    queues, hash tables, and heaps will all be built on top of
    \texttt{std::vector} or the indexing pattern it exposes. The
    habits you lock in here -- correct sizing, bounds checking,
    pass-by-reference, reserve-before-loop -- will carry through
    every remaining chapter.
\end{itemize}

\begin{notebox}[What this connects to]
\begin{itemize}
    \item \textbf{Chapter 2} takes the loops and iteration patterns you
    learned here and reframes them as \emph{problem-solving strategies}
    (recursion, greedy, DP). The loops don't change; the way you decide
    \emph{which loop to write} does.
    \item \textbf{Chapter 3} formalizes the cost analysis you've been doing
    intuitively (``\texttt{push\_back} is cheap, \texttt{insert} at the front
    is not''). Big-O is the notation; the underlying mechanisms -- shifting,
    reallocation, amortization -- are exactly what you practiced here.
    \item \textbf{Chapter 4} introduces \emph{linked lists} as the alternative
    to contiguous storage. Every performance argument in that chapter is
    phrased as ``compared to a vector, $\ldots$'' -- so the more solid your
    vector intuitions are here, the faster ch.~4 lands.
    \item \textbf{Chapter 5+} (stacks, queues, deques) are built on top of
    vector and linked-list primitives. The \texttt{push\_back/pop\_back}
    pattern from 1.7 \emph{is} the stack interface.
\end{itemize}
\end{notebox}

\textbf{Companion materials.} A two-page compact notes reference for this
chapter lives at \texttt{chapters/ch\_1/notes.tex}; twelve drill prompts you
can paste into a fresh Claude session live at
\texttt{chapters/ch\_1/practice.md}. Use the notes to review, the prompts to
test.

\end{document}
